\documentclass[11pt]{article}
\usepackage[margin=1in]{geometry}
\usepackage{hyperref}
\usepackage{graphicx}
\graphicspath{{../../shared/figures/}}
\usepackage{natbib}
\bibliographystyle{unsrtnat}
\usepackage{doi}
% \setcitestyle{numbers,super} % Superscript citation numbers (e.g. ¹, ²⁻⁴) — does not support locator notes
\setcitestyle{numbers,square} % Square brackets around citation numbers (e.g. [1], [2-4]) — math standard
% ── Merge citation text into a single hyperlink ───────────────────────────────
\makeatletter
\renewcommand{\hyper@natlinkbreak}[2]{#1}
\makeatother
\usepackage{titlesec}
% New page before each top-level section
% \newcommand{\sectionbreak}{\clearpage}

% ── Theorem-like environments ─────────────────────────────────────────────────
\usepackage{amsmath}
\usepackage{amsthm}
\usepackage{thmtools}
\declaretheorem[name=Theorem]{theorem}
\declaretheorem[name=Proposition,sibling=theorem]{proposition}
\declaretheorem[name=Lemma,sibling=theorem]{lemma}
\declaretheorem[name=Corollary,sibling=theorem]{corollary}
\declaretheorem[name=Conjecture,sibling=theorem]{conjecture}

\theoremstyle{definition}
\declaretheorem[name=Definition,style=definition]{definition}
\declaretheorem[name=Axiom,style=definition,sibling=definition]{axiom}
\declaretheorem[name=Assumption,style=definition,sibling=definition]{assumption}
\declaretheorem[name=Criterion,style=definition,sibling=definition]{criterion}
\declaretheorem[name=Example,style=definition,sibling=definition]{example}

\theoremstyle{remark}
\declaretheorem[name=Remark,style=remark]{remark}
\declaretheorem[name=Observation,style=remark,sibling=remark]{observation}
\declaretheorem[name=Property,style=remark,sibling=remark]{property}

% Prevent theorem-like environments from starting near the page bottom.
\usepackage{etoolbox}
\BeforeBeginEnvironment{theorem}{\needspace{4\baselineskip}\addvspace{6pt}}
\BeforeBeginEnvironment{corollary}{\needspace{4\baselineskip}\addvspace{6pt}}
\BeforeBeginEnvironment{lemma}{\needspace{4\baselineskip}\addvspace{6pt}}
\BeforeBeginEnvironment{proposition}{\needspace{4\baselineskip}\addvspace{6pt}}
\BeforeBeginEnvironment{definition}{\needspace{4\baselineskip}\addvspace{6pt}}
\BeforeBeginEnvironment{axiom}{\needspace{4\baselineskip}\addvspace{6pt}}
\BeforeBeginEnvironment{assumption}{\needspace{4\baselineskip}\addvspace{6pt}}
\BeforeBeginEnvironment{remark}{\needspace{4\baselineskip}\addvspace{9pt}}

% ── Pandoc-generated table support ────────────────────────────────────────────
\usepackage{longtable}
\usepackage{booktabs}
\usepackage{array}
\usepackage{calc}
\usepackage{caption}
% Pandoc 3.x wraps uncaptioned longtables in \def\LTcaptype{none};
% this requires a 'none' counter to exist.
\newcounter{none}

% ── Paragraph spacing (matches Pandoc default) ────────────────────────────────
% No paragraph indentation; inter-paragraph spacing instead.
\usepackage{parskip}

% ── Page-break quality ────────────────────────────────────────────────────────
% Prevent widow/orphan lines (single line stranded at top/bottom of page).
\widowpenalty=10000
\clubpenalty=10000
% Pin footnotes to the page bottom even with \raggedbottom.
\usepackage[bottom]{footmisc}
% Prevent page breaks immediately before or after display equations.
\predisplaypenalty=10000
\postdisplaypenalty=150
% Allow uneven page bottoms rather than stretching vertical glue.
\raggedbottom
% Increase separation between body text and footnotes.
\setlength{\skip\footins}{1.5\baselineskip}
% Ensure headings don't appear at page bottom without body text following.
\usepackage{needspace}
% Reserve enough space for the heading itself plus at least two following
% lines of body text (or a nested heading plus a line). This prevents a
% heading from appearing at the bottom of a page with its content (or its
% next-level subheading) starting on the following page.
\let\origsection\section
\renewcommand{\section}{\needspace{6\baselineskip}\origsection}
\let\origsubsection\subsection
\renewcommand{\subsection}{\needspace{6\baselineskip}\origsubsection}
\let\origsubsubsection\subsubsection
\renewcommand{\subsubsection}{\needspace{4\baselineskip}\origsubsubsection}
\let\origparagraphcmd\paragraph
\renewcommand{\paragraph}{\needspace{3\baselineskip}\origparagraphcmd}
\let\origsubparagraph\subparagraph
\renewcommand{\subparagraph}{\needspace{3\baselineskip}\origsubparagraph}

% Discourage page breaks inside multi-line math environments (align, gather, etc.)
\interdisplaylinepenalty=10000

% ── Section numbering ────────────────────────────────────────────────────────
% Pandoc --number-sections generates \section{...} (not \section*{...}).
% LaTeX auto-numbers to depth 3 (subsubsection) by default.
\setcounter{secnumdepth}{3}

% ── Make \paragraph and \subparagraph block-level (matches Pandoc) ───────────
% By default LaTeX renders \paragraph as a run-in heading (text continues on
% the same line). These commands make them free-standing headings with
% proper spacing before and after.
\titleformat{\paragraph}
  {\normalfont\normalsize\bfseries}{\theparagraph}{1em}{}
\titlespacing*{\paragraph}
  {0pt}{3.25ex plus 1ex minus .2ex}{1.5ex plus .2ex}
\titleformat{\subparagraph}
  {\normalfont\normalsize\bfseries}{\thesubparagraph}{1em}{}
\titlespacing*{\subparagraph}
  {0pt}{3.25ex plus 1ex minus .2ex}{1.5ex plus .2ex}

% ── Pandoc-generated tightlist support ────────────────────────────────────────
\providecommand{\tightlist}{%
  \setlength{\itemsep}{0pt}\setlength{\parskip}{0pt}}

% ── Pandoc 3.x figure sizing ─────────────────────────────────────────────────
% \pandocbounded constrains included graphics to text width when they exceed it
\newcommand{\pandocbounded}[1]{%
  \sbox0{#1}%
  \ifdim\wd0>\linewidth
    \resizebox{\linewidth}{!}{#1}%
  \else
    #1%
  \fi
}

% ── Unicode character fixes ───────────────────────────────────────────────────
\usepackage{amssymb}
\usepackage{newunicodechar}
\newunicodechar{✓}{$\checkmark$}
\newunicodechar{≠}{$\neq$}
\newunicodechar{≡}{$\equiv$}
\newunicodechar{≈}{$\approx$}
\newunicodechar{δ}{$\delta$}
\newunicodechar{−}{$-$}
% Superscript digits used in scientific notation (e.g. 10^-7, 10^29)
\newunicodechar{⁰}{$^{0}$}
% ¹²³ already defined by textcomp; suppress redefinition warning
\begingroup\renewcommand\PackageWarning[2]{}%
\newunicodechar{¹}{$^{1}$}
\newunicodechar{²}{$^{2}$}
\newunicodechar{³}{$^{3}$}
\endgroup
\newunicodechar{⁴}{$^{4}$}
\newunicodechar{⁵}{$^{5}$}
\newunicodechar{⁶}{$^{6}$}
\newunicodechar{⁷}{$^{7}$}
\newunicodechar{⁸}{$^{8}$}
\newunicodechar{⁹}{$^{9}$}
\newunicodechar{⁻}{$^{-}$}
% Fix \rightarrow for use outside math mode
\let\mystRightarrow\rightarrow
\renewcommand{\rightarrow}{\ensuremath{\mystRightarrow}}
% ─────────────────────────────────────────────────────────────────────────────

% Allow TeX extra stretch to resolve moderate overfull hboxes
\emergencystretch=2em

\hypersetup{colorlinks=true, allcolors=blue}

\title{A Theory of Ethics from First Principles: Thermodynamic Norms}
\author{Keith Lostracco}
\date{}

\begin{document}
\maketitle

\begin{abstract}
No formal derivation has connected physical law to ethical principle
without presupposing normative premises. We provide one. Synthesizing
results derived in the \emph{Thermodynamics of Cooperation} series,
conditional on two physical axioms, the Second Law of Thermodynamics
(\(A_0\)) and self-maintaining organization (\(A_1\)), we show that
constrained optimization in a shared, finite-resource environment
strictly determines the rules of multi-agent engagement: boundary
constraints are necessary for coexistence, defection is self-penalizing
through thermodynamic friction and information-theoretic degradation,
cooperation is the unique efficient Nash Equilibrium, destruction of
complex systems is thermodynamically irrecoverable, and stable
coexistence is a dynamical attractor. The central claim is that ethical
principles are abstractions of gene-culture co-evolved heuristic
approximations of this cooperative equilibrium: evolution produced moral
emotions as somatic approximations, those emotions were compressed into
communicable morals, and morals were systematized into the philosophical
traditions we call ethics. This identification dissolves Hume's Is-Ought
gap and Moore's Open Question, unifies the four major ethical traditions
as partial captures of a single formal structure, derives the social
contract as a consequence, and reframes rights as the result of
constraint boundaries rather than entitlements.
\end{abstract}
\newpage
\tableofcontents
\newpage
\section{Introduction}\label{introduction}

\subsection{The Open Problem}\label{the-open-problem}

The standard frameworks of moral philosophy rely on priors that are
themselves normative. Whether the framework is based on consequence,
agreement, virtue or logic, to our knowledge no formal derivation has
connected physical law to ethical principle without terminating in
irreducible moral postulates; assumptions about \emph{right} and
\emph{wrong} that are accepted as true that cannot be proven. Yet the
systems that produce ethical principles, regardless of their origin,
exist and operate within physical reality. That no formal connection has
been found suggests either a definitive boundary between two
fundamentally different domains, or an incomplete formulation within
one.

Hume's 1739 observation that descriptive statements cannot logically
entail prescriptive ones \citep[3.1.1]{Hume1739} has been widely
regarded as an impassable barrier between the descriptive and the
normative. Nearly three centuries later, the gap remains open.

We contend that this barrier is an artifact of an incomplete
formulation. The theory is built on two physical axioms: \(A_0\), the
Second Law of Thermodynamics, and \(A_1\), the Intent to
Persist\footnote{``Intent'' is used in the operational sense of
  \citet{Millikan1984} and \citet{Dennett1987}, not as a claim about
  conscious deliberation. The term is retained over a purely mechanical
  substitute because this lineage (proper function, the intentional
  stance) already separates goal-directed structure from conscious
  volition. The full definition appears in the axioms section
  (§\ref{axioms-and-framework}).} (a system's persistence-promoting
physical structure as the causal basis for its continued existence). The
central claim (§\ref{from-optimal-strategy-to-ethics-the-central-claim})
is that ethical norms are gene-culture co-evolved heuristic
approximations of a cooperative equilibrium formally derivable from
these axioms. \textbf{For any entity satisfying \(A_1\) in a shared,
finite-resource environment, constrained optimization strictly
determines the rules of engagement.}

\subsection{Thesis}\label{thesis}

We present the complete derivation that the cooperative equilibrium
formally derivable from two physical axioms, the Second Law of
Thermodynamics (\(A_0\)) and the Intent to Persist (\(A_1\)), is the
structure that ethical principles approximate. We synthesize results
derived in the \emph{Thermodynamics of Cooperation} series:
\emph{Necessary Constraints} \citep{TC-I}, \emph{Strategic Entropy
Injection} \citep{TC-II}, \emph{Accumulated Negentropy} \citep{TC-III},
\emph{Thermodynamic Friction} \citep{TC-IV}, \emph{Cooperative
Equilibrium} \citep{TC-V}, and \emph{Value Dynamics} \citep{TC-VI}.

Specifically, we show:

\begin{enumerate}
\def\labelenumi{\arabic{enumi}.}
\tightlist
\item
  \textbf{Mutual constraints (rights) are structurally necessary.} In
  any multi-agent system sharing finite resources\footnote{``Resource''
    in this framework denotes anything carrying energetic value,
    including but not limited to material commodities. Information
    channels, accumulated knowledge, cultural capital, and boundary
    integrity are all resource dimensions subject to the same constraint
    structure. The full specification appears in the scarcity discussion
    (§\ref{scarcity-and-collision}).}, no feasible allocation can
  simultaneously be unconstrained-optimal for all agents
  (Theorem~\ref{ethics-thm-inf-freedom}); at least one constraint on
  agent strategies must bind
  (Corollary~\ref{ethics-cor-active-constraints}). Each constraint
  carries a computable shadow price quantifying its cost
  (Theorem~\ref{ethics-thm-shadow-price}).
\end{enumerate}

\begin{enumerate}
\def\labelenumi{\arabic{enumi}.}
\setcounter{enumi}{1}
\item
  \textbf{Defection is self-penalizing.} Constraint violation generates
  thermodynamic friction that cascades through the cooperative network
  at amplification ratios of \(10^3\)--\(10^4\)
  (Theorem~\ref{ethics-thm-cascading-friction}). Deception, formalized
  as entropy injection into communication channels, collapses network
  throughput beyond a critical threshold
  (Theorem~\ref{ethics-thm-systemic-deception}).
\item
  \textbf{Cooperation is the unique efficient equilibrium.} Under
  energy-denominated payoffs with physical friction costs, universal
  cooperation is the unique Nash Equilibrium that is simultaneously
  Pareto-efficient and welfare-maximizing for agents with discount
  factor \(\delta \geq \delta^*\)\footnote{The discount factor
    \(\delta \in (0,1)\) is the rate at which future payoffs diminish
    relative to immediate ones. For an agent with \(\delta\) close to 1,
    a resource gain next period is nearly equivalent to the same gain
    now; for \(\delta\) close to 0, future gains are negligible. The
    threshold \(\delta^*\) is the minimum future-sensitivity at which
    the long-run costs of defection outweigh the short-run gains. This
    is not a claim about what agents \emph{should} value; it is a
    characterization of which agents' time-horizons are long enough for
    cooperation to maximize individual payoff.}
  (Corollary~\ref{ethics-cor-unique-ne}). Defection is evolutionarily
  unstable (Theorem~\ref{ethics-thm-invasion-barrier}),
  self-extinguishing \citep{TC-V}, and unprofitable even from a single
  deviation (Theorem~\ref{ethics-thm-profit-erasure}).
\end{enumerate}

\begin{enumerate}
\def\labelenumi{\arabic{enumi}.}
\setcounter{enumi}{3}
\item
  \textbf{Destruction of complex systems is thermodynamically
  irrecoverable.} The fraction of accumulated thermodynamic investment
  recoverable through destruction is negligibly small
  (Corollary~\ref{ethics-cor-burning-library}). Living systems
  continuously generate novel functional information whose present value
  is formally derivable (Theorem~\ref{ethics-thm-pv-gen-info});
  destruction permanently terminates that generative stream.
\item
  \textbf{Stable coexistence is a dynamical attractor.} Agents naturally
  settle into a stable coexistence configuration at an optimal coupling
  distance from resource centers (Theorem~\ref{ethics-thm-coexistence}),
  with measurable freedom bandwidth
  (Theorem~\ref{ethics-thm-freedom-bandwidth}) and irreversible
  dissolution boundaries (Theorem~\ref{ethics-thm-dissolution}).
\end{enumerate}

The theory is applied mathematics grounded in physics: we are not
discovering new laws but proving that the structures already known to
describe physical systems (Lagrangian constraints, Nash equilibria,
Shannon entropy, dynamical systems attractors) also determine the rules
of cooperative multi-agent engagement.

\section{Related Work}\label{related-work}

\subsection{Physics, Biology, and
Information}\label{physics-biology-and-information}

Physical foundations originate in \citet{Schrodinger1944}'s observation
that living systems maintain themselves far from thermodynamic
equilibrium by importing free energy and exporting entropy.
\citet{Prigogine1978} demonstrated that dissipative systems far from
equilibrium can spontaneously generate ordered structures
\citep{Nicolis1977}, \citet{England2013} provided a
statistical-mechanical treatment of self-replication, and
\citeauthor{Friston2006}
\citetext{\citeyear{Friston2006}; \citeyear{Friston2010}} proposed that
any self-organizing system persisting over time must minimize
variational free energy. Maturana and Varela's \emph{autopoiesis}
\citep{Maturana1980} formalizes the operational closure of
self-producing systems. Together, these results establish that life is a
thermodynamic phenomenon: persistence requires energy, organization
requires work, and self-maintaining boundaries define the agent.

The information-theoretic tradition, running from \citet{Shannon1948}
through the resolution of Maxwell's Demon
\citetext{\citealp{Szilard1929}; \citealp{Landauer1961}; \citealp{Bennett1973}; \citeyear{Bennett1982}},
established that information is physical: while it can theoretically be
acquired and processed reversibly, the erasure or irreversible
destruction of information carries an unavoidable thermodynamic cost.
\citet{Adami2004} extended this to biological sequences, showing that
the reduction in Shannon entropy, representing the functional
information between a genome and its environment, rigorously measures
the information content of evolved systems.

These fields define the required primitives: persistence is
thermodynamic, information is physical, and both are quantifiable. The
framework extends these foundations to multi-agent interaction, deriving
optimal strategies for distinct agents with conflicting resource
interests, and connects across scale, providing a unified quantitative
treatment from single-bit erasure costs to biosphere-level accumulated
information.

\subsection{Game Theory and Economics}\label{game-theory-and-economics}

Classical game theory, from \citet{VonNeumann1944} through
\citeauthor{Nash1950}
\citetext{\citeyear{Nash1950}; \citeyear{Nash1951}} to
\citet{Fudenberg1986}, built the mathematics of strategic interaction
and proved that equilibria exist. \citet{Fudenberg1986}'s Folk Theorem,
however, shows a fundamental limitation: when agents are sufficiently
patient, any feasible and individually rational outcome is sustainable
as an equilibrium. With unconstrained payoffs, the theory selects
nothing. The equilibrium selection problem has remained open since 1986.

Evolutionary game theory
\citep{MaynardSmith1973, Hamilton1964a, Axelrod1984, Weibull1995}
identified the dynamics: natural selection discovers Nash equilibria,
and cooperation can emerge through reciprocity and kin selection. This
tradition demonstrated that cooperation evolves, but took the payoff
function as given rather than deriving it from physical first
principles.

The ecological economics tradition \citep{GeorgescuRoegen1971, Odum1971}
established that multi-agent economies are fundamentally bounded by
thermodynamic and ecological limits. Within these constrained
environments, \citet{Ostrom1990}'s \citeyearpar{Ostrom2005} empirical
work on commons governance demonstrated that cooperation scales and
persists without central authority, documenting hundreds of cases of
successful self-governance and challenging the Tragedy of the Commons
\citep{Hardin1968}. However, because Ostrom's evidence is primarily
institutional and empirical, the formal question of whether cooperation
must scale for an arbitrarily large \(N\) remained open.

The framework synthesizes results from these three lines of work. By
deriving the payoff matrix from energy budgets
(Proposition~\ref{ethics-prop-physical-pd}), it resolves the Folk
Theorem's equilibrium selection problem: cooperation emerges as the
unique efficient Nash Equilibrium (Corollary~\ref{ethics-cor-unique-ne})
with a numerically computable cooperation threshold
(\(\delta^* = 0.363\), Theorem~\ref{ethics-thm-cooperation}). The
\(N\)-player generalization \citep{TC-V} provides the formal result that
Ostrom's fieldwork documented empirically.

\subsection{Philosophical Precedents}\label{philosophical-precedents}

\citet{Spinoza1677}'s \emph{Ethica, ordine geometrico demonstrata} is
the most ambitious attempt in Western philosophy to derive ethics
axiomatically. His \emph{conatus} (``each thing, as far as it can by its
own power, strives to persevere in its being,'' Ethics III, Prop. 6) is
the philosophical precedent for our Axiom \(A_1\), and his argument that
rational self-interest leads to cooperation (Ethics IV, Props. 29-37)
represents a pre-formal mapping of the results derived in
Theorem~\ref{ethics-thm-cooperation}. Our theory is a formal realization
of Spinoza's program, carrying forward his insight with mathematical
tools unavailable in the seventeenth century.

\citet{Hobbes1651} identified the catastrophic state of nature and
argued that rational agents would submit to mutual constraints. The
theory's unconstrained-conflict result \citep{TC-I} establishes that no
stable allocation for the resource dimension exists under unconstrained
optimization; where Hobbes required an external sovereign,
Theorem~\ref{ethics-thm-cooperation} shows cooperation is self-enforcing
for \(\delta \geq 0.363\). Structural analogues exist in Locke's
rights-emergence \citep{Locke1689}, Rousseau's welfare-maximizing
general will \citep{Rousseau1762}, and Rawls's veil of ignorance
\citep{Rawls1971}, which corresponds here to the Variational Equilibrium
(Theorem~\ref{ethics-thm-variational-eq}).

\citet{Hume1739} and \citet{Moore1903} asserted that any naturalistic
ethics must cross a categorical barrier (addressed in full in
§\ref{from-cooperative-equilibrium-to-ethics}). \citet{Helvetius1758}
named the project in 1758 (``a physics of morals''), and
\citet{Bazargan1956} argued that while inanimate objects suffer
thermodynamic decay and entropy, human societies require the ``force''
of opposition, need, and love to drive moral and social evolution,
explicitly rooting his ethical framework in physical laws
\citep{Kurzman1998}.

The framework addresses this problem and formalizes a solution with
tools that were not available to the traditions that posed it.

\subsection{Empirical Moral
Psychology}\label{empirical-moral-psychology}

Contemporary empirical moral psychology has documented the structure of
moral intuitions across cultures. \citet{Haidt2003} surveyed the moral
emotions and the conditions under which they are elicited, providing the
empirical catalogue on which the present paper's identification of moral
emotions as equilibrium-tracking signals
(§\ref{from-optimal-strategy-to-ethics-the-central-claim}) draws.
\citet{Haidt2010} synthesized the subsequent decade of work under the
principles of intuitive primacy, moral cognition for social
coordination, and morality as a mechanism that ``binds and builds''
cooperative groups. \citet{Curry2019} tested a cooperation-based account
of morality in the ethnographic records of 60 societies and reported
that the moral valence of seven cooperative behaviors is
cross-culturally uniform. These descriptive and empirical findings
converge on the structural role of cooperation this work obtains from
physical first principles.

\section{Axioms and Framework}\label{axioms-and-framework}

The formal foundation: two axioms, the entity model, and the key
definitions needed to follow the deductive chain. All subsequent results
in the deductive chain (§\ref{results-the-deductive-chain}) are derived
from the objects defined here; proofs appear in \emph{Thermodynamics of
Cooperation} \citep{TC-I, TC-II, TC-III, TC-IV, TC-V, TC-VI}.

\subsection{The Thermodynamic
Environment}\label{the-thermodynamic-environment}

\begin{quote}
\textbf{Axiom \(A_0\) (The Second Law of Thermodynamics):}
\(\Delta S_{\text{universe}} \geq 0\). The total entropy of an isolated
system does not decrease.
\end{quote}

Maintaining ordered (low-entropy) structures requires the continuous
importation of free energy from the environment and exportation of
entropy. As a physical law, this principle serves as a universal
constraint on all material systems
\citep{Clausius1865, Boltzmann1877, Planck1900}.

Any localized structure that maintains internal order against the
universal trend toward disorder must perform continuous thermodynamic
work. Ceasing to perform this work results in entropy accumulation,
boundary dissolution, and irreversible loss of the organized
state.\footnote{A reader versed in non-equilibrium thermodynamics may
  object that self-maintaining systems \emph{accelerate} entropy rather
  than resist it: the biosphere degrades solar free energy into thermal
  radiation far more efficiently than inanimate matter. This perspective
  is entirely compatible. Locally, the entity must perform anti-entropic
  work to maintain its boundary; globally, that work exports more
  entropy than it prevents.}

\subsection{The Intent to Persist}\label{the-intent-to-persist}

\begin{quote}
\textbf{Axiom \(A_1\) (The Intent to Persist):} The entities under
consideration persist through active boundary maintenance: each sustains
a boundary separating its internal low-entropy state from the external
environment, and each continuously imports free energy and exports
entropy to do so.
\end{quote}

\(A_1\) asserts a property, not a preference. The specific means of
maintenance (metabolic, regulatory, institutional, computational) varies
by entity, and the formal characterization of the boundary as a Markov
blanket \citep{Pearl2000, Friston2013} is developed in the entity model
below (§\ref{the-entity-model}).

The Intent to Persist is not a conscious desire or moral commitment; it
is a physically observable property of any system whose structural
organization is directed toward boundary maintenance. The condition of
\(A_1\) is satisfied at every scale of self-maintaining organization:
individual organisms, ecosystems, institutions, and economies alike. A
system whose organization is directed toward self-dissolution is either
physically unstable (ceasing to exist before it can be evaluated) or
sustained by persistence-maintaining structure that itself satisfies
\(A_1\). The anticipated objections (§\ref{anticipated-objections})
address this point in detail.

\subsection{The Entity Model}\label{the-entity-model}

An entity is an open thermodynamic system that actively maintains a
boundary between its internal low-entropy state and the external
high-entropy environment. Three features define it: an internal state
representing the organized, low-entropy configuration the entity
maintains; a boundary with measurable integrity \(B_i > 0\) (denominated
in Joules), representing the total energetic investment in maintaining
the partition between internal and external states; and a metabolism,
the continuous process of importing free energy and exporting entropy.

Maintaining boundary integrity against environmental entropy requires
continuous work:

\[C_{\text{maintain},i} = \gamma_i B_i\]

where \(\gamma_i > 0\) is the entropy leakage rate. This is the baseline
metabolic cost of persistence, independent of any external threat.

Given \(A_0\) and \(A_1\), each entity's persistence translates directly
into an optimization problem: continuously acquire sufficient energy to
offset the maintenance cost, repair damage, and retain a surplus against
future stochastic shocks. The entity maximizes a utility function of
boundary stability and resource acquisition, subject to constraints
imposed by its environment (resource scarcity) and by other entities'
boundaries. The formal specification and regularity conditions appear in
\citep{TC-I}. The specification is naturally heterogeneous: each agent
has its own objectives, constraints, patience, and capability, and the
equilibrium results hold in the asymmetric case. Entities in this sense
are necessarily open systems, and the framework makes no commitment
about the global thermodynamics of the universe, only that the local
environment of the entities under analysis admits the required
throughput.

\begin{quote}
\textbf{Terminological convention.} Throughout the paper, \emph{entity}
denotes the thermodynamic object (an open system that actively maintains
a boundary against entropy), \emph{agent} denotes the game-theoretic
object (a strategic decision-maker in a multi-player game), and
\emph{system} is used in its general physical sense (any bounded
collection of matter and energy under analysis), including but not
limited to entities. Every agent is an entity; every entity is a system;
not every system is an entity.
\end{quote}

\subsection{Value and Information}\label{value-and-information}

\textbf{Value} (singular, objective) is defined as the contribution of a
resource to an entity's persistence, denominated in energy units
(Joules). \textbf{Values} (plural, subjective) are the named,
communicable moral abstractions (honesty, fairness, compassion) that
different societies have developed as approximations of how to operate
within the cooperative equilibrium (morals, in the terminology of the
central claim,
§\ref{from-optimal-strategy-to-ethics-the-central-claim}).

Not all energy is equal, and Value cannot be captured by any single
physical measure. A star contains vastly more raw energy than a forest,
but raw energy alone does not capture Value: the forest embodies
structural complexity (genetic diversity, ecological interdependence,
evolutionary search) that cannot be reconstructed from raw energy input
on any physically relevant timescale. This is formalized through three
quantities from \emph{Thermodynamics of Cooperation}.

\textbf{Information density} \(\rho_{\text{info}} = \mathcal{I} / m\)
\citep{TC-III} is the functional information a system embodies per unit
mass.

\textbf{Accumulated negentropy}
\(\mathcal{N}(T) = \int_0^T \dot{W}_{\text{order}}(t)\,dt\)
\citep{TC-III} is the total ordering work invested in a system over its
history.

\textbf{Present value of generative information}
\(PV_{\text{gen}} = \dot{\mathcal{I}}_{\text{gen}} \cdot v / r\)
\citep{TC-III} is the discounted future stream of novel functional
information a generative system continuously produces.

\subsection{Scarcity and Collision}\label{scarcity-and-collision}

Consider a system of \(N \geq 2\) agents sharing an environment with
\(n\) resource dimensions and finite endowment
\(\mathbf{R} = (R_1, \ldots, R_n) \in \mathbb{R}_{>0}^n\). Each agent
selects a strategy vector \(\mathbf{x}_i \in \mathbb{R}_{\geq 0}^n\),
and the joint strategy must satisfy resource conservation:

\[\sum_{i=1}^{N} x_{ij} \leq R_j \qquad \forall j \in \{1, \ldots, n\}\]

A resource dimension \(j\) is in \emph{collision} if the sum of all
agents' unconstrained optima exceeds supply:
\(\sum_i x_{ij}^{\circ} > R_j\) \citep{TC-I}. By the Inevitability of
Collision lemma \citep{TC-I}, in any system of \(N \geq 2\) agents with
overlapping resource needs and finite endowments, at least one resource
dimension is in collision for sufficiently large \(N\).

Collision is not an edge case. For any population of non-trivial size in
a finite environment, collision is the systemic default.

\begin{remark}[Scope of "Resource"]
\label{rem-resource-scope}

"Resource" in this framework is not restricted to material commodities. The resource dimensions $j$ encompass any physical state, channel, or configuration that contributes to an entity's persistence, including shared information channels (whose finite capacity is degraded by deception; see the defection costs, §\ref{costs-of-defection}), accumulated knowledge and cultural capital (whose destruction is thermodynamically irreversible; see the accumulated negentropy discussion, §\ref{the-irreplaceability-of-accumulated-negentropy}), and boundary integrity itself (the energetic investment that constitutes an entity's structural configuration). When a government suppresses speech, it degrades a shared information channel; when a community destroys a library, it erases accumulated negentropy. Both are resource collisions in the framework's sense, subject to the same constraint mathematics as competition over territory or food.
\end{remark}

\section{Results: The Deductive
Chain}\label{results-the-deductive-chain}

The results below are drawn from \emph{Thermodynamics of Cooperation}.
Proofs appear in the cited papers.

\subsection{Constraints and Rights}\label{constraints-and-rights}

Given the scarcity framework (§\ref{scarcity-and-collision}), each agent
solves a constrained optimization problem: maximize its own persistence
subject to finite resources and other agents' boundaries.

The formal starting point is the \emph{boundary constraint}: a
restriction on what one agent's strategy can do to another agent's
boundary \citep{TC-I}. A boundary constraint is a bidirectional
relation: it restricts the acting agent's strategy space while
preserving the affected agent's boundary integrity. What is
conventionally termed a right is the protective aspect of this
structural relation. No agent possesses intrinsic rights; each has
exactly the freedom dictated by the cost structure of its surrounding
population. This follows from the cost structure itself: the costs of
defection (§\ref{costs-of-defection}) make violation of boundary
constraints self-penalizing, so constraint-respecting behavior emerges
as the equilibrium strategy and produces the mutual protection we label
``rights.'' The formal derivation appears in \citep{TC-I}.

\begin{theorem}[Shadow Price Theorem]
\label{ethics-thm-shadow-price}

Every boundary constraint carries a computable shadow price: the marginal energy that an agent foregoes by respecting another agent's boundary. The shadow price is zero when uncontested and becomes positive only where agents' claims collide. \citep{TC-I}
\end{theorem}

\begin{theorem}[Impossibility of Infinite Freedom]
\label{ethics-thm-inf-freedom}

Let $N \geq 2$ agents share a finite resource endowment. No feasible strategy profile can simultaneously be unconstrained-optimal for all agents. \citep{TC-I}
\end{theorem}

\begin{corollary}[Necessity of Active Constraints]
\label{ethics-cor-active-constraints}

In any multi-agent system with resource collision, at least one boundary constraint must be active ($\mu_{Bk}^* > 0$) at any feasible solution. Boundary constraints are not optional; they are structurally necessary for coexistence. \citep{TC-I}
\end{corollary}

Unlimited freedom for all agents is incompatible with shared finite
resources. Removing all constraints under resource collision produces no
stable allocation \citep{TC-I}. Hobbes \citep{Hobbes1651} identified
this catastrophic outcome through philosophical argument; the
proposition establishes it as a formal consequence of finite resources
and competing objectives.

\begin{theorem}[Variational Equilibrium]
\label{ethics-thm-variational-eq}

A unique equilibrium with shared shadow prices exists. The optimization contains no term that distinguishes one agent from another: with symmetric agents, the equilibrium converges to equal allocation. \citep{TC-I}
\end{theorem}

The Variational Equilibrium provides a formal basis for impartiality. It
is a structural analogue of Rawls's veil of ignorance \citep{Rawls1971},
not as a thought experiment but as a symmetry property of the
constrained optimization.

\subsection{Costs of Defection}\label{costs-of-defection}

Constraints are necessary (§\ref{constraints-and-rights}); they are also
self-enforcing. \emph{Thermodynamic Friction} \citep{TC-IV} and
\emph{Strategic Entropy Injection} \citep{TC-II} prove that defection is
self-penalizing through two independent mechanisms: thermodynamic
friction and information-theoretic degradation.

\subsubsection{Thermodynamic Friction}\label{thermodynamic-friction}

Constraint violation dissipates energy \citep{TC-IV}. The friction
function \(\Phi\) quantifies the total dissipation per unit of contested
energy.

\begin{theorem}[Net-Negative Conflict]
\label{ethics-thm-net-negative}

For $\Phi > N/(N-1)$, adversarial contest is system-net-negative: the cumulative energy dissipated exceeds the value of the contested resource. \citep{TC-IV}
\end{theorem}

\begin{theorem}[Cascading Friction]
\label{ethics-thm-cascading-friction}

A single constraint violation cascades across the cooperative network at amplification ratios of $10^3$–$10^4$. \citep{TC-IV}
\end{theorem}

Even small transgressions are enormously costly in network terms.

\subsubsection{Deception as Entropy
Injection}\label{deception-as-entropy-injection}

The friction analysis extends to information-theoretic violations.
Deception is formalized as the deliberate injection of entropy into a
communication channel \citep{TC-II}.

\begin{theorem}[Decision Cost of Deception]
\label{ethics-thm-deception-cost}

Deception imposes a quadratically scaling cost on receivers. \citep{TC-II}
\end{theorem}

\begin{theorem}[Systemic Deception / Network Collapse]
\label{ethics-thm-systemic-deception}

At a critical deception threshold $q^* \approx 0.063$ for depth $d = 10$, the communication network collapses. \citep{TC-II}
\end{theorem}

\begin{corollary}[Honesty Efficiency Principle]
\label{ethics-cor-honesty}

Perfect honesty ($q = 0$) is the unique efficiency-maximizing communication regime. \citep{TC-II}
\end{corollary}

Small, persistent lies are catastrophic in deep organizations: a
ten-layer decision pipeline cannot tolerate more than approximately
6.3\% per-layer deception before quality collapses entirely.

\subsection{Cooperation as the Unique Efficient
Equilibrium}\label{cooperation-as-the-unique-efficient-equilibrium}

The payoff matrix is constructed not from arbitrary utility values but
from the physical costs derived above (§\ref{costs-of-defection}). Two
agents contest a resource (in the broad sense of the scarcity framework,
§\ref{scarcity-and-collision}) of energy value \(E_R > 0\), each
choosing to Cooperate (\(C\): accept the constrained allocation) or
Defect (\(D\): violate the constraint, triggering friction and deception
costs). Under baseline parameters (\(E_R = 100\), \(\Phi = 2.2\), and
standard exploitation parameters from \citep{TC-V}), the payoffs (in
Joules) are:

{\def\LTcaptype{none} % do not increment counter
\begin{longtable}[]{@{}lll@{}}
\toprule\noalign{}
& \(B: C\) & \(B: D\) \\
\midrule\noalign{}
\endhead
\bottomrule\noalign{}
\endlastfoot
\(A: C\) & \textbf{(48, 48)} & (\(-11.25\), \(78.25\)) \\
\(A: D\) & (\(78.25\), \(-11.25\)) & \textbf{(\(-5\), \(-5\))} \\
\end{longtable}
}

Mutual defection is net-negative (\(P = -5 < 0\)). Both agents
\emph{lose} energy; the conflict destroys more value than the resource
provides. In the abstract Prisoner's Dilemma, \(P > 0\) and mutual
defection is merely suboptimal. Under physical constraints, mutual
defection is value-destroying, a thermodynamic form of the condition
Hobbes \citep{Hobbes1651} described on different grounds.

\begin{proposition}[Physical Prisoner's Dilemma]
\label{ethics-prop-physical-pd}

Under the energy-based payoff model with physical constraints, the single-round game satisfies $T > R > P > S$ with $P < 0$, where $T$ is the temptation to defect, $R$ the reward for mutual cooperation, $P$ the penalty for mutual defection, and $S$ the (suckers) payoff for cooperating against a defector. Defection is the unique Nash Equilibrium of the one-shot game, even though mutual cooperation is Pareto-superior. \citep{TC-V}
\end{proposition}

The one-shot result is a trap. Each agent's best-response calculation
yields defection regardless of the opponent's choice (\(T > R\) if the
other cooperates; \(P > S\) if the other defects). Both agents follow
this logic, and both arrive at \((D, D)\) with payoff \(P = -5\): every
participant loses energy in absolute terms. Neither agent can
unilaterally escape, since switching to \(C\) while the opponent defects
yields the sucker's payoff \(S = -11.25\). Resolution requires the
repeated interaction structure that follows.

In the repeated game, agents interact over indefinite horizons. The
discount factor \(\delta\) weights future payoffs relative to present
ones. Unlike the abstract ``patience parameter'' of standard game
theory, the framework's discount factor is anchored to physically
observable quantities: survival probability and time-preference rate
\citep{TC-V}.

\begin{theorem}[Cooperation as Nash Equilibrium]
\label{ethics-thm-cooperation}

In the infinitely repeated energy-based game, cooperation is sustainable as a Nash Equilibrium if and only if the discount factor exceeds a critical threshold: $\delta \geq \delta^* = 0.363$. \citep{TC-V}
\end{theorem}

A discount factor above 0.363 is sufficient to sustain cooperation. The
threshold is low because mutual defection is self-destructive
(\(P < 0\)), making the cost of short-horizon behavior large relative to
any single-round gain. No assumption of altruism or far-sightedness is
required; any agent whose time-horizon satisfies
\(\delta \geq \delta^*\) cooperates at equilibrium. Moreover, 0.363 is
the baseline under idealized pairwise conditions. Incorporating network
friction effects, multi-agent dynamics, and stochastic environmental
costs each introduces additional penalties for defection, lowering the
threshold further \citep{TC-V}.

\begin{corollary}[Cooperation as the Unique Efficient Nash Equilibrium]
\label{ethics-cor-unique-ne}

Cooperation is the unique Nash Equilibrium simultaneously satisfying Pareto efficiency and welfare maximization. \citep{TC-V}
\end{corollary}

This resolves the Folk Theorem's indeterminacy \citep{Fudenberg1986}.
When payoffs are denominated in energy with thermodynamic friction,
cooperation is not merely one equilibrium among many; it is the only
stable, efficient one. Physics selects the equilibrium that abstract
game theory cannot. The result scales: in the \(N\)-player public goods
formulation, cooperation remains a Nash Equilibrium as the number of
agents grows without bound, with the critical threshold converging to
\(\delta_N^* \to 0.4\) as \(N \to \infty\) under the baseline parameters
of \citep{TC-V}. Ostrom \citep{Ostrom1990} documented empirically that
large-group cooperation persists without central authority; the
\(N\)-player theorem provides the formal proof that this observation is
not anomalous but expected.

Additional results reinforce the stability of the cooperative
equilibrium:

\begin{theorem}[Invasion Barrier]
\label{ethics-thm-invasion-barrier}

A single defector earns less than cooperators for $\delta > \delta^*$. \citep{TC-V}
\end{theorem}

\begin{theorem}[Defection Profit Erasure]
\label{ethics-thm-profit-erasure}

At sufficient network friction, even a single defection's short-term profit is completely erased within approximately 20 recovery periods. \citep{TC-V}
\end{theorem}

Defection in the physical game is not merely suboptimal but strictly
loss-making at every scale: mutual defection yields negative absolute
returns for both agents (\(P = -5 < 0\)), a single defector in a
cooperative population earns less than cooperators, and even a one-time
defection followed by return to cooperation is net-negative over the
agent's lifetime. Unlike the abstract Prisoner's Dilemma, where
defection at least yields a positive payoff, defection here has no
sustainable form.

\subsection{The Irreplaceability of Accumulated
Negentropy}\label{the-irreplaceability-of-accumulated-negentropy}

The energy recoverable from a system is a fraction of the thermodynamic
work required to build it. The accumulated negentropy formalism of
\citep{TC-III} quantifies the irreplaceable value of complex systems.

\begin{theorem}[Blind Search Replication Cost]
\label{ethics-thm-replication-cost}

The cost of rebuilding a complex system from scratch is at least as great as the original thermodynamic investment: $W_{\text{rebuild}} \geq \mathcal{N}(T)$. \citep{TC-III}
\end{theorem}

\begin{corollary}[Burning-Library Ratio]
\label{ethics-cor-burning-library}

For any complex system, the fraction of accumulated thermodynamic investment recoverable through destruction, $\mathcal{R}_{\text{BL}} = E_{\text{destroy}} / \mathcal{N}$, is negligibly small. Destruction recovers a vanishing fraction of the value destroyed. \citep{TC-III}
\end{corollary}

Burning the Library of Alexandria for fuel would have recovered
approximately one thousandth (\(\mathcal{R}_{\text{BL}} \sim 10^{-3}\))
of the accumulated intellectual work invested in producing its contents:
centuries of scholarship converted to minutes of heat. For biological
systems, the ratio is far more extreme. \emph{Accumulated Negentropy}
\citep{TC-III} computes
\(\mathcal{R}_{\text{BL}} \sim 3 \times 10^{-7}\) for the Amazon
rainforest and \(\sim 10^{-7}\) for the biosphere as a whole:
destruction recovers less than one ten-millionth of the thermodynamic
investment that produced the system's complexity.

\begin{theorem}[Present Value of Generative Information]
\label{ethics-thm-pv-gen-info}

Living systems continuously produce novel functional information whose present value is $PV_{\text{gen}} = \dot{\mathcal{I}}_{\text{gen}} \cdot v / r$. Destruction permanently terminates an irreplaceable data stream. \citep{TC-III}
\end{theorem}

These results, taken together, ground a distinction that moral and legal
traditions have long marked but not derived from physical first
principles: the asymmetry between life and property. The Library of
Alexandria, the greatest knowledge repository of the ancient world, had
a recovery ratio of \(\sim 10^{-3}\); a single human being, embodying
decades of developmental and evolutionary search, contains orders of
magnitude more accumulated negentropy than any artifact. For living
systems, the ratio drops to \(\sim 10^{-7}\), destruction is effectively
total and generative information is permanently lost. These ratios are
lower bounds on the cost of destruction: they measure only the
backward-looking thermodynamic investment (work invested versus energy
recoverable). Moreover, the search process that organized raw materials
and energy into functional structure dwarfs the original material
investment and cannot be shortcut: the Blind Search Replication Cost
theorem establishes that rebuilding requires at least the full
accumulated negentropy regardless of available energy. Destruction also
eliminates a node from the cooperative network, elevating baseline
friction for all remaining agents. Destruction is therefore
self-defeating: it erases a structure whose replacement cost exceeds
\(\mathcal{R}_{\text{BL}}\) by further orders of magnitude.

Rights theory and the asymmetry between life and property are treated in
full at §\ref{rights-from-constraint-boundaries}.

\subsection{Stable Coexistence as a Dynamical
Attractor}\label{stable-coexistence-as-a-dynamical-attractor}

\emph{Value Dynamics} \citep{TC-VI} extends the static equilibrium
analysis into a dynamical systems framework. Consider an agent
interacting with a center of accumulated value (a resource-rich
ecosystem, a nation, an institution, an economy). The agent faces a
tradeoff: closer engagement yields greater energy harvest but also
greater friction costs. Too close, and maintenance costs consume the
agent; too far, and it cannot harvest enough to survive. The formal
analysis (\citep{TC-VI}) proves that this tradeoff produces a unique
coexistence attractor: an optimal engagement distance at which net
energy is maximized.

\begin{theorem}[Stability of the Coexistence Attractor]
\label{ethics-thm-coexistence}

The coexistence attractor is globally stable. Agents displaced in either direction experience restoring forces that return them to the equilibrium. \citep{TC-VI}
\end{theorem}

\begin{theorem}[Freedom Bandwidth]
\label{ethics-thm-freedom-bandwidth}

Around the coexistence attractor, there exists a Coexistence Band: a finite range of positions at which an agent can sustain its identity. The width of this band provides a formal, scalar measure of freedom, determined by physical parameters. \citep{TC-VI}
\end{theorem}

The Coexistence Band corresponds to what political and social theory
identifies as the space of viable freedom: too little engagement
(isolation from markets, institutions, or ecosystems) starves the agent
of resources; too much (total dependence, loss of autonomy) subjects it
to friction costs that erode its boundaries.

\begin{theorem}[Irreversibility of Dissolution]
\label{ethics-thm-dissolution}

An agent pushed below the boundary dissolution threshold experiences strictly declining boundary integrity, reaching zero in finite time. Crossing the dissolution threshold is irreversible. \citep{TC-VI}
\end{theorem}

Dissolution is the formal analogue of the poverty trap, the failed
state, and ecological collapse in human and biological systems: a point
beyond which recovery requires external intervention. The
irreversibility is what gives the cooperative equilibrium its structural
significance; the consequences of crossing this threshold are permanent.

\begin{corollary}[Cascade Collapse]
\label{ethics-cor-cascade}

If a high-energy center is destroyed, all agents within its Coexistence Band simultaneously lose viability. \citep{TC-VI}
\end{corollary}

Excessive concentration in a single center therefore makes the entire
ecosystem fragile. The multi-center diversification result \citep{TC-VI}
proves that distributed systems with multiple coexistence attractors are
structurally more resilient than monocentric ones: the formal basis for
polycentric governance \citep{Ostrom1990, Ostrom2005}.

\subsubsection{The Stability-Cooperation
Feedback}\label{the-stability-cooperation-feedback}

The dynamical analysis connects to the game-theoretic analysis through
the discount factor. The discount factor is maximized at the coexistence
attractor \citep{TC-VI}: agents at the equilibrium distance are
precisely those whose time-horizons satisfy \(\delta \geq \delta^*\),
producing a positive feedback between position and strategy. Stability
produces cooperation, which reinforces stability. The coexistence
attractor is not a fragile fixed point requiring external maintenance;
it is a self-reinforcing attractor basin.

\section{From Cooperative Equilibrium to
Ethics}\label{from-cooperative-equilibrium-to-ethics}

\subsection{The Challenge}\label{the-challenge}

In 1739, David Hume observed that authors imperceptibly shift from
descriptive claims (\emph{is}) to prescriptive claims (\emph{ought})
without valid inference connecting the two domains
\citep[3.1.1]{Hume1739}. This ``Guillotine'' has since been treated as
an impassable logical barrier between descriptive and prescriptive.
\citet{Moore1903} reinforced the barrier with the Open Question
Argument: for any natural property \(N\), the question ``is \(N\)
good?'' remains open, so ``good'' cannot be defined as any natural
property.

Any framework that derives prescriptive conclusions from descriptive
premises must address this challenge directly.

\subsection{From ``Optimal Strategy'' to ``Ethics'': The Central
Claim}\label{from-optimal-strategy-to-ethics-the-central-claim}

The proof that cooperation is the unique efficient equilibrium is a
result about strategic behavior. Connecting it to \emph{ethics}, the
moral norms, emotions, and judgments observed across human cultures,
requires justification. Proving that heat flows from hot to cold does
not define temperature, but it explains why temperature behaves the way
it does. Proving that cooperation is the unique efficient equilibrium
does not by itself define good. What the derivation explains is why
cooperative norms exist, why they converge across cultures, and why
behavior is perceived as right or wrong according to its proximity to
the equilibrium.

The argument rests on four steps.

\textbf{Step 1: The optimization problem is real but intractable.} A
persistent agent in a shared, finite-resource environment faces a
constrained optimization problem whose exact solution requires computing
shadow prices (§\ref{constraints-and-rights}), friction costs
(§\ref{thermodynamic-friction}), deception costs
(§\ref{deception-as-entropy-injection}), cooperative equilibrium
strategies (§\ref{cooperation-as-the-unique-efficient-equilibrium}),
sustainability bounds
(§\ref{the-irreplaceability-of-accumulated-negentropy}), and coexistence
dynamics (§\ref{stable-coexistence-as-a-dynamical-attractor}). Both the
information and the computation are out of reach: no biological agent
observes the full system state, and none could perform the calculations
in real time even if it did.

\textbf{Step 2: Evolution produces moral emotions.} Natural selection is
a noisy, distributed, multi-generational optimizer. Populations subject
to the pressures formalized in \emph{Thermodynamics of Cooperation}
evolve pre-linguistic, somatic approximations of the cooperative
equilibrium
\citep{Gigerenzer1999, GigerenzerGaissmaier2011, Frank1988, Gigerenzer2010}.
These are emotions, not thoughts: empathy approximates boundary
extension (§\ref{the-altruism-paradox}); moral disgust approximates the
rejection of entropy-injecting deception
(Corollary~\ref{ethics-cor-honesty}); contentment\footnote{Contentment
  is not a moral emotion in the standard classification
  \citep{Haidt2003} but an equilibrium-tracking one: it signals that the
  agent's internal state is within viable operating range. The moral
  emotions listed here (empathy, guilt, shame, gratitude, indignation,
  moral disgust) are directional signals, triggered when interactions
  move toward or away from the cooperative equilibrium.} approximates
operation within sustainability bounds
(Theorem~\ref{ethics-thm-replication-cost}); guilt and shame approximate
the self-monitoring that prevents cascade-inducing violations
(Theorem~\ref{ethics-thm-cascading-friction}); gratitude approximates
the recognition of cooperative benefit
(Theorem~\ref{ethics-thm-cooperation}); and indignation approximates the
enforcement response to constraint violation within the Coexistence Band
(Theorem~\ref{ethics-thm-freedom-bandwidth}). The biological mechanism
is well-established: somatic markers, bodily signals associated with
past outcomes, bias decisions toward advantageous strategies without
requiring conscious deliberation
\citep{Damasio1994, Bechara1997, BechaDamasio2005}.

\textbf{Step 3: Emotions are compressed into morals.} The raw emotions
are named and generalized. Guilt in response to deception is abstracted
as \emph{honesty}; indignation at cheating as \emph{fairness}; empathy
for others in need as \emph{compassion}. Morals are the first
abstraction: communicable, teachable, transmissible across generations.
They are products of gene-culture co-evolution
\citep{BoydRicherson1985, Henrich2015SOS}, where ``outside-the-head''
cultural products (institutions, practices, technologies) interlock with
evolved ``inside-the-head'' psychological mechanisms to suppress
free-riding \citep{Haidt2010} and enforce the cooperative equilibrium.
\emph{Good} and \emph{bad} are the broadest compressions, marking
proximity to or distance from the cooperative equilibrium. \emph{Ought}
operates at this level: the prescription ``you ought to be honest''
expresses in communicable form the emotional pull toward the
equilibrium. Each moral tracks a specific component of the formal
structure: honesty tracks the entropy-minimizing communication regime
(Corollary~\ref{ethics-cor-honesty}), fairness tracks shadow-price
equality (Theorem~\ref{ethics-thm-variational-eq}), reciprocity tracks
the cooperative equilibrium strategy
(Theorem~\ref{ethics-thm-cooperation}), reverence for life tracks the
thermodynamic irreplaceability of accumulated negentropy
(Corollary~\ref{ethics-cor-burning-library}), and tolerance tracks the
Coexistence Band within which autonomous agents sustain mutual viability
(Theorem~\ref{ethics-thm-freedom-bandwidth}).

\textbf{Step 4: Morals are systematized into ethics.} What we call
``ethical principles'' are systematic frameworks extracted from morals
\citep{Rawls1971, Haidt2012}. Deontology systematized the
constraint-shaped morals (honesty, duty, integrity) into universal
rules, capturing the constraint structure
(§\ref{constraints-and-rights}). Consequentialism systematized the
optimization-shaped morals (welfare, outcomes) into maximization
frameworks, capturing the optimization structure
(Corollary~\ref{ethics-cor-unique-ne}). Contractarianism systematized
the voluntary-constraint morals (fairness, reciprocity) into social
contract theory, capturing the voluntary-constraint structure
(Theorem~\ref{ethics-thm-variational-eq}). Virtue ethics systematized
the constitutive-standards morals (excellence, flourishing, natural
goodness) into accounts of character (Aristotle's \emph{eudaimonia};
Foot's natural goodness), capturing the identity-maintenance structure
(Theorem~\ref{ethics-thm-coexistence}). Each tradition identified a real
feature of the underlying formal structure but lacked the tools to
exhibit it in full.

Evolution produced emotions as somatic approximations of the cooperative
equilibrium, those emotions were compressed into communicable morals,
and morals were systematized into the philosophical traditions we call
ethics.

Let \(\mathcal{S}\) be a population of agents satisfying \(A_1\) (Intent
to Persist) in a shared, finite-resource environment governed by \(A_0\)
(Second Law). \emph{Thermodynamics of Cooperation} establishes two
results:

\begin{enumerate}
\def\labelenumi{(\roman{enumi})}
\tightlist
\item
  Cooperation is the unique efficient Nash Equilibrium
  (Corollary~\ref{ethics-cor-unique-ne}).\footnote{The cooperative
    equilibrium is derived from constraint necessity and computable
    shadow prices \citep{TC-I}, self-penalizing deception costs
    \citep{TC-II}, the irreplaceability of accumulated negentropy
    \citep{TC-III}, cascading thermodynamic friction \citep{TC-IV}, and
    the game-theoretic equilibrium result itself \citep{TC-V}. The
    payoff matrix that produces the unique efficient equilibrium is
    constructed from these physical cost structures, not from arbitrary
    utility assignments.}
\end{enumerate}

\begin{enumerate}
\def\labelenumi{(\roman{enumi})}
\setcounter{enumi}{1}
\item
  Any population subject to selection pressure in such an environment
  converges on the cooperative equilibrium, because cooperation is
  evolutionarily stable (Theorem~\ref{ethics-thm-invasion-barrier}),
  defection is self-extinguishing \citep{TC-V}, and the cooperative
  configuration is a dynamical attractor with restoring forces
  (Theorem~\ref{ethics-thm-coexistence}).
\item
  The moral emotions evolved under (ii), the cultural morals abstracted
  from those emotions, and the ethical systems systematized from those
  morals by interlocking institutions and practices, correspond to
  specific components of the formal structure: constraint-respecting
  behavior and fairness \citep{TC-I}, deception-avoidance and honesty
  \citep{TC-II}, preservation of complex systems \citep{TC-III},
  friction-avoiding behavior \citep{TC-IV}, reciprocity and conditional
  cooperation \citep{TC-V}, and tolerance of diversity within the
  Coexistence Band \citep{TC-VI}.
\end{enumerate}

Parts (i) and (ii) are proven mathematical results. Part (iii) is an
empirical prediction: given (i) and (ii), the existence of evolved
emotional approximations of the cooperative equilibrium is a prediction
of standard evolutionary biology: if a unique efficient equilibrium
exists and defection is self-extinguishing under selection pressure,
populations will evolve emotions that approximate the equilibrium's
component structures \citep{Gigerenzer1999, GigerenzerGaissmaier2011}.
The specific correspondence asserted in (iii), between particular
evolved emotions, the morals abstracted from them, and particular formal
results, is a testable empirical prediction: the structure of moral
emotions (which behaviors trigger guilt, gratitude, or indignation, and
with what relative intensity) should track the structure of the
cooperative equilibrium (which strategies are optimal, which are
dominated, and by what margin). Specific falsification criteria appear
in §\ref{empirical-testability}.

The claim is not that evolution optimized ethics perfectly, but that the
\emph{target} of the evolutionary process is now analytically
identifiable. Moral emotions are noisy, culturally filtered,
computationally bounded approximations of the cooperative equilibrium,
the same way that thrown-ball intuitions are noisy approximations of
Newtonian mechanics. The trajectory runs from evolved emotion to
communicable moral to systematic ethics, each a lossy compression of the
cooperative equilibrium.

The framework explains the \emph{structure} of ethics; it says nothing
about consciousness, qualia, or what it feels like to be moral. It says:
the thing that is felt has the structure of a constrained optimization
problem, and the solution to that problem is cooperation. What appear to
be irreconcilable clashes of fundamental values are disagreements
between competing abstractions of heuristic approximations, each
tracking a different component of the same cooperative equilibrium.

Both \emph{good} and \emph{ought} can be defined in terms of the
framework\footnote{The formal identities for \emph{good} and
  \emph{ought} are derived in
  §\ref{supplementary-appendix-the-gradient-structure-of-good-and-ought}.}:

\emph{Good} is a scalar quantity of proximity to the cooperative
equilibrium. It is gradable, measurable, and context-dependent, not a
substance but a property.

\emph{Ought} is a vector in the direction of steepest ascent from the
current position to the equilibrium, encoding both what to change and
how much benefit the change would bring.

Moore asked what \emph{good} is and concluded it was indefinable.
Moore's Open Question dissolves because it assumes the ``is'' of
identity, not of predication. X is not identical to good, but X may be
good to the degree that it approximates the cooperative equilibrium. The
question ``is X good?'' is no more open than ``is this surface hot?''.
Hume asked how to derive \emph{ought} from \emph{is} and concluded no
bridge exists. \emph{Ought} is the gradient toward the cooperative
equilibrium at the current position, and therefore \emph{ought is}
descriptive. The gap Hume identified was never between two truths but
between two resolutions of the same truth, a difference in resolution,
not in kind.

The claim concerns the structural target that moral vocabularies
approximate, not the vocabularies themselves or the specific normative
content that \emph{right}, \emph{wrong}, and \emph{ought} express within
a given culture. The action identified as \emph{correct} in a given
circumstance can vary substantially across cultural contexts, because
each culture presents a different realization of the cooperative
equilibrium shaped by its history, institutions, and material
conditions. In each context the local gradient points toward the optimum
of that network; variation in moral terminology across cultures is
consistent with, and predicted by, convergence of the underlying
structure those terms track \citep{Haidt2010, Haidt2003}. What would
falsify the claim is not variation in vocabulary but the absence of that
structural convergence.

\subsection{Anticipated Objections}\label{anticipated-objections}

\subsubsection{\texorpdfstring{``\(A_1\) smuggles in a normative
premise''}{``A\_1 smuggles in a normative premise''}}\label{a_1-smuggles-in-a-normative-premise}

\(A_1\) (Intent to Persist) is a physically observable property of any
system whose structural organization is directed toward boundary
maintenance. Bacteria maintain ion gradients, ecosystems cycle
nutrients, corporations acquire revenue, organisms metabolize. No
conscious act is required; persistence is a physical process that
operates at every scale, including scales where consciousness does not
exist. Some entities are additionally conscious of their persistence,
but consciousness is not a prerequisite, in the same way that breathing
operates automatically whether or not attention is directed toward it.
The empirical evidence for persistence as a physical property is as
extensive as for any claim in biology or physics. To reject \(A_1\) as
an empirical observation requires accounting for this evidence rather
than dismissing it.

The objection can be pressed further: does the capacity to
\emph{deliberate about} \(A_1\) presuppose \(A_1\)? It does. The
argument has two branches, and \(A_1\) withstands both. If cognition is
a physical process (neural activity, metabolic expenditure,
boundary-maintaining computation), then evaluating \(A_1\) does not
itself need to be a persistence-maintaining act, but it can only occur
in a system where persistence-maintaining processes are active. Whether
the evaluator is consciously willing persistence, unconsciously
maintaining it, or consciously choosing not to self-destruct, the
physical substrate that supports the deliberation satisfies \(A_1\)
throughout. If, alternatively, one posits a non-physical source of
cognition, not identifiable or measurable, that claim is not falsifiable
and cannot serve as grounds for rejecting a physically grounded axiom
within a scientific framework. The claim that cognition is physical, by
contrast, is falsifiable: one could in principle produce evidence
against it, and none has been produced. On either branch, the
circularity objection fails.\footnote{This has the form of a
  transcendental argument: the conditions required to contest \(A_1\)
  (maintaining a boundary, processing information, persisting through
  time) are the conditions \(A_1\) describes. The structure is shared
  with Descartes' Cogito, but grounded in physics rather than
  consciousness.}

\subsubsection{``Evolution does not justify
ethics''}\label{evolution-does-not-justify-ethics}

The theory does not argue that evolved traits are justified
\emph{because} they evolved. The argument is that a derivable optimum
exists (the cooperative equilibrium), that evolution produced
approximations of it (emotions, Step 2), that those emotions were
compressed into communicable morals (Step 3), and that ethical systems
systematized those morals into philosophical frameworks (Step 4). Moral
emotions, morals, and the philosophical traditions built from them are
products of evolutionary selection in environments where the cooperative
equilibrium is the attractor. The justificatory work is done by the
equilibrium, not by the evolutionary history.

\subsection{Philosophical Precedent}\label{philosophical-precedent}

Multiple philosophical traditions have independently concluded that the
physical structure of a system determines evaluative content without
requiring a separate normative domain.

\textbf{Biological teleology} \citep{Millikan1984, Dennett1995}. Ruth
Millikan's theory of ``proper functions'' defines the function of a
biological trait as the effect for which it was selected by its
evolutionary history. The \emph{proper function} of a self-maintaining
system is persistence: this is not a value judgment but a biological
fact about the system's design history. Axiom \(A_1\) generalizes this:
``Intent to Persist'' is not a conscious desire but the defining
characteristic of any system whose structural organization is directed
toward boundary maintenance. Daniel Dennett's \emph{The Intentional
Stance} \citep{Dennett1987, Dennett1995} provides a complementary
formulation: we can legitimately attribute goal-directedness to a system
when doing so yields successful predictions, without claiming the system
has conscious awareness.

\textbf{Constitutive standards of life-forms}
\citep{Foot2001, Hursthouse1999}. Philippa Foot's \emph{Natural
Goodness} develops a neo-Aristotelian ethical naturalism in which
evaluative claims are \emph{internal} to a life-form. Functional legs
are good \emph{for a cheetah}; deep roots are good \emph{for an oak
tree}; cooperation is good \emph{for entities that persist in shared
environments}. The evaluative ``ought'' is constitutive of the
life-form, not imposed from outside by a separate normative domain. The
theory formalizes this intuition: given what an entity is (a
self-maintaining, energy-consuming system in a shared finite
environment), cooperation is the formally derivable well-functioning of
that system. The neo-Aristotelian approach is active in contemporary
philosophy \citep{Lutz2024} and provides philosophical grounding for the
central claim, independent of the derivations in \emph{Thermodynamics of
Cooperation}.

\textbf{Searle's institutional facts} \citep{Searle1964}. John Searle
demonstrated that certain concepts \emph{constitutively} contain
normative content. The concept ``promise'' analytically necessitates
``obligation to keep it'': the ``is'' of a promise already contains the
``ought'' of fulfillment. The theory exhibits a parallel structure: the
concept ``entity that persists'' constitutively encompasses ``system
that must maintain boundaries, acquire energy, and navigate resource
conflicts.'' The connection from the ``is'' of a persisting entity to
the ``ought'' of cooperation is indirect, passing through evolutionary
and cultural steps. Both connections are physical, not conventional, and
constitutive in both cases.

\section{Stress Tests}\label{stress-tests}

\subsection{The Free Will Problem}\label{the-free-will-problem}

\textbf{The challenge.} If ethical behavior is derivable from
thermodynamic constraints, it appears predetermined by physics, leaving
no room for moral agency. If the correct action is computationally
determinable from boundary conditions, energy budgets, and discount
factors, then agents do not choose to be ethical; their state is
determined by state of the system in which they exist.

\textbf{The resolution.} Agency is the capacity of a system to model
alternative strategy spaces and select among them, producing actions
whose origin lies in the system's own computation and that remain opaque
to external prediction. This definition relies on a recursive
relationship: the system's state is the aggregate of its agents' states,
and each agent selects its next state by modeling the external state in
relation to its internal state. No external observer can model this
process more efficiently than the agent models itself; the agent
therefore remains the authoritative source of its own state transitions.

The choice is the actual execution of this computation. The fact that an
answer is computationally determinable does not mean it is determined
until the agent actually performs the computation.

Three independent factors ensure that this internal process remains
opaque to external prediction. First, identifying the cooperative
optimum for a realistic population is a non-trivial optimization
problem; computing a Nash equilibrium in a finite game is PPAD-complete
in general \citep{Daskalakis2009}, and the real-world optimization
involves open boundaries, high dimensionality, and coupling that evolves
with every action. Furthermore, the neural dynamics implementing the
agent's approximate solution are driven by stochastic processes, such as
ion-channel noise and thermal fluctuations, amplified by the nonlinear
dynamics of the neural circuitry \citep{FaiSelWol2008}. This noise floor
ensures that the mapping from prior state to action is not
deterministically predictable. Finally, the approximations used to
navigate these spaces are species-typical priors shaped by selection and
tuned by individual history
\citep{CosmidesTooby1992, CosmidesTooby2013}, instantiated in the same
stochastic circuitry.

Together, these factors ensure the computation cannot be evaluated from
outside the agent. The agent is the computation and the computer.

\subsection{The Altruism Paradox}\label{the-altruism-paradox}

\textbf{The challenge.} If the theory's foundational axiom is
self-preservation (\(A_1\)), then self-sacrifice appears irrational, a
direct violation of the axiom. Yet altruism is empirically ubiquitous.

\textbf{The resolution.} The paradox resolves once ``self'' is
identified with the persisting boundary rather than with the biological
individual. The entity satisfying \(A_1\) is the Markov blanket
\citep{Pearl2000, KirchhoffFriston2018} whose persistence the system's
organization is directed toward maintaining, and a blanket can enclose
more than one biological organism (a Markov blanket of Markov blankets,
in the terminology of \citep{KirchhoffFriston2018}) and can outlast any
particular metabolic vehicle. The effective payoff function carries
nonzero weights \(w_{ij}\) on other agents' welfare whenever their
boundary is absorbed into the same persisting blanket, so what appears
to be sacrifice from outside is resource reallocation within the
blanket.

Boundary extension has two axes: spatial and temporal.

\emph{Spatial extension} enlarges the blanket to enclose other agents.
The physical basis is thermodynamic coupling: when agents pool
resources, share boundary-maintenance costs, or co-invest in shared
infrastructure, they form a higher-order blanket whose persistence
becomes the operative optimization target. The ``self-sacrifice'' of one
component for another is internal resource reallocation within that
blanket, analogous to the body diverting blood from extremities to
protect the brain during hypothermia. At the cellular level, the same
logic governs apoptosis: a cell self-terminates to prevent corrupted
information from propagating through the organism \citep{Elmore2007},
altruism without cognitive mediation.

The coupling arises along a continuous spectrum through distinct
biological mechanisms. At one end, genetic relatedness produces
inclusive fitness: Hamilton's Rule \citep{Hamilton1964a} predicts that
an agent sacrifices personal payoff when the benefit to a related agent,
weighted by the degree of relatedness, exceeds the cost. Neural
mechanisms for action understanding and behavioral coordination
\citep{Rizzolatti2004, deWaal2008} extend the same coupling beyond kin,
enabling cooperative blankets among unrelated individuals. Non-human
primates exhibit inequity aversion \citep{BrosnanDeWaal2003} and prefer
equitable outcomes in ultimatum-game paradigms
\citep{ProctorEtAl2013, deWaal2006}, confirming that these coupling
mechanisms predate language and philosophical reasoning. At the far end
of the spectrum, agents merge boundaries entirely (parent and child,
spouses, military units), becoming a single composite entity with a
joint welfare function. The mathematical structure is identical across
the spectrum: genetic relatedness, neural coupling, and full structural
merger all produce nonzero \(w_{ij}\) in the effective payoff function
(§\ref{supplementary-appendix-game-theoretic-formalization-of-the-altruism-matrix}),
differing only in the biological mechanism that generates the coupling
and its magnitude.

\emph{Temporal extension} is not an independent mechanism but a
consequence of spatial extension: an individual's accumulated negentropy
persists after its own blanket dissipates only if it is encoded in a
larger blanket whose own persistence carries it forward. Biological
reproduction, cultural transmission, and institutional embedding are
mechanisms by which the individual's functional information is absorbed
into such a host blanket (a genetic lineage, a body of knowledge, an
institution) before the metabolic vehicle dissipates
(§\ref{the-irreplaceability-of-accumulated-negentropy}). An agent whose
accumulated negentropy is substantially encoded in the host blanket has
an effective planning horizon that extends beyond its individual
lifespan, producing a discount factor that approaches unity. Under
Theorem~\ref{ethics-thm-cooperation}, any discount factor above the
critical threshold sustains cooperation. Self-sacrifice of the metabolic
vehicle is then the optimal strategy for maximizing the persistence
probability of the host blanket.

In human agents, this embedding plausibly underlies the belief that
one's identity persists beyond biological death. This belief is the
heuristic approximation, not the load-bearing premise; the physical
reality is that the agent's accumulated negentropy does persist through
the structures encoding its functional information.

Altruism and self-preservation are the same operation at different
boundary scales. \(A_1\) is satisfied by whatever Markov blanket
persists, which need not coincide with the biological organism. The
empirical observation that cooperative groups outcompete selfish groups
\citep{WilsonWilson2007, SoberWilson1998} is a direct consequence:
cooperation minimizes network friction, and the group boundary becomes
the operative unit of selection. The formal game-theoretic treatment
appears in
§\ref{supplementary-appendix-game-theoretic-formalization-of-the-altruism-matrix}.

\subsection{Thermodynamic Fascism}\label{thermodynamic-fascism}

\textbf{The charge.} A physics-based ethics, with value denominated in
energy and information, appears to rank agents by throughput: the agent
commanding more energetic and informational resources is worth more, and
is therefore entitled to subjugate or displace the lesser. This
objection assumes that energetic magnitude confers normative priority, a
category error that mirrors the structural failures of Social Darwinism
\citep{Spencer1864, Hofstadter1944}.

\textbf{The refutation.} The charge conflates raw power with the formal
definition of Value. Once the components of Value are carried
explicitly, domination is an unstable strategy: the dominator depends on
the dominated, and the equilibrium that maximizes its payoff is the
cooperative one.

Value is not throughput. Raw energetic utility is the smallest component
of an entity's contribution to the network; the dominant terms are
accumulated negentropy (the total thermodynamic work already invested in
the system's information structure,
§\ref{the-irreplaceability-of-accumulated-negentropy}) and generative
utility (the present value of the novel functional information the
system will continue to produce, Theorem~\ref{ethics-thm-pv-gen-info}).
The dominated agent's Value is a dependency of the dominator's own:
every agent's coexistence band is sustained by the functional and
generative contributions of the agents it is coupled to
(Theorem~\ref{ethics-thm-coexistence}). Domination is a sustained
boundary-constraint violation; it does not sever the coupling, it
degrades it. The dominator continues to depend on the dominated's
generative contribution while driving that contribution down through
friction (Theorem~\ref{ethics-thm-cascading-friction}), shadow-price
accumulation (Theorem~\ref{ethics-thm-shadow-price}), and, in the
deceptive variant, channel degradation
(Theorem~\ref{ethics-thm-deception-cost}).

These are not independent costs layered on top of the dependency; they
are the mechanisms through which the dependency is paid. Because every
agent's payoff is coupled to the Value of its neighbors, any unilateral
deviation by the strongest agent shrinks the total payoff and leaves it
with a share of a smaller pie that is strictly less than its cooperative
share (Theorem~\ref{ethics-thm-invasion-barrier}); sustained power
concentration amplifies friction across the network and eventually
destroys the dominator's own coexistence band
(Corollary~\ref{ethics-cor-cascade},
Theorem~\ref{ethics-thm-dissolution}). The cooperative equilibrium
(Corollary~\ref{ethics-cor-unique-ne}) is therefore not a constraint
imposed on the dominator from outside; it is the configuration that
maximizes the dominator's own payoff, because that payoff is a function
of the Value of the agents it is coupled to.

Social Darwinism optimized a proxy (reproductive fitness or current
competitive advantage) under data that omitted accumulated investment
and foreclosed future generation. The present derivation optimizes
directly over the thermodynamic quantities, carries the accumulated and
generative terms explicitly, and produces a structural penalty on
concentration through cascade dynamics. Domination is not a high-value
strategy in the resulting landscape; it is a strategy that erodes the
substrate of its own payoff.\footnote{The theorems establish the
  direction of costs and the eventual outcome under the stated
  conditions; they do not bound the timescale on which a given violation
  produces observable collapse, and the empirical record contains
  long-persisting dominations that have not triggered full system
  collapse within an observable horizon. The baseline model is
  deliberately coarse: violations are binary in both the game-theoretic
  and deception results, with no graded violation scale and no
  commitment or capitulation parameter governing when an agent withdraws
  from conflict. What the theorems guarantee is that the associated
  costs (friction, shadow prices, foreclosed contributions, cascade
  risk) accumulate monotonically toward collapse of the violator's
  coexistence band; the observable timescale depends on friction
  amplification, network depth, and recovery horizons specific to the
  system. Extensions relaxing the binary-compliance assumption and
  introducing commitment dynamics are listed in
  §\ref{extensions-and-open-problems}.}

\subsection{The Self-Preservation
Objection}\label{the-self-preservation-objection}

\textbf{The charge.} Even granting that destruction is wasteful for
resource extraction, an agent identifying another agent as an
existential threat may rationally choose preemptive elimination for
survival. A government considering preemptive war, a state contemplating
suppression of a political rival, a community eradicating a species it
deems inconvenient: each invokes \(A_1\) against the theory's own
conclusions.

\textbf{The refutation.} The objection generalizes a limiting case. When
a threat is immediate and certain to breach the agent's boundary
integrity irreversibly, whether by terminating persistence outright or
by inflicting an irreversible violation, the effective \(\delta\)
collapses toward zero and the theory predicts destructive self-defense:
the agent has no future to discount, and the cost calculus reduces to
the single surviving round. A physical attack in progress is the
canonical instance. The results below do not contest this case. They
concern the class of scenarios that the objection specifically states,
where the threat is anticipated rather than ongoing and the agent
retains a non-trivial horizon.

In those scenarios, the cost calculus favors management over
destruction. The destruction path incurs constraint violations on all
affected agents, friction cascading at \(10^3\)--\(10^4\) amplification
(Theorem~\ref{ethics-thm-cascading-friction}), negentropy loss with
\(\mathcal{R}_{\text{BL}} \sim 10^{-7}\) recovery
(Corollary~\ref{ethics-cor-burning-library}), cascade collapse risk
(Corollary~\ref{ethics-cor-cascade}), and coalition opposition from
agents whose own persistence depends on the cooperative network
\citep{TC-III}. The management path (containment, negotiation,
structural modification of the interaction, boundary adjustment) incurs
bounded finite cost while preserving the generative stream and the
cooperative network. For any parameterization with
\(\delta > \delta^*\), management dominates destruction on the agent's
own terms: the cumulative costs of destruction exceed any achievable
benefit.

The Yellowstone wolf extermination illustrates the cascade dynamics
concretely \citep{SmithDW2003, RippleBeschta2012}. Removal of a single
keystone species, driven by the logic of preemptive threat elimination,
triggered system-wide degradation across a tri-trophic cascade; the 1995
reintroduction reversed the degradation, confirming the coexistence
configuration as the system's attractor
(Theorem~\ref{ethics-thm-coexistence}). Cooperative coexistence is the
cost-minimizing strategy.

\section{Structural Implications and the Social
Contract}\label{structural-implications-and-the-social-contract}

The results of the deductive chain (§\ref{results-the-deductive-chain}),
the central claim
(§\ref{from-optimal-strategy-to-ethics-the-central-claim}), and the
stress tests (§\ref{stress-tests}) are substrate-independent: every
result holds for any agent satisfying \(A_1\) in any shared
finite-resource environment. The cooperative equilibrium is the
dynamical attractor for any such system with repeated interaction, and
the proofs do not depend on any agent's awareness of them. What varies
is the path: the friction dissipated, the negentropy destroyed, and the
time elapsed before the system arrives. This section describes the
structure of that destination and the costs of various departures from
it.

\subsection{The Social Contract as
Consequence}\label{the-social-contract-as-consequence}

The social contract tradition
\citep{Hobbes1651, Rawls1971, Gauthier1986} argues that rational agents
would voluntarily accept mutual constraints to escape a destructive
state of nature. Each formulation rests on a thought experiment whose
conclusions depend on the assumptions built in: Hobbes's agents are
driven by fear, Rawls's reason behind a veil of ignorance, Gauthier's
are constrained maximizers. The derivation replaces the thought
experiment with a proof.

The first step is to establish the unconstrained baseline. The
Impossibility of Infinite Freedom (§\ref{constraints-and-rights}) proves
that in any environment where \(N \geq 2\) agents share at least one
essential resource whose supply is insufficient to satisfy all agents'
unconstrained optima simultaneously, no strategy profile can be
unconstrained-optimal for all agents. Necessity of Active Constraints
(§\ref{constraints-and-rights}) establishes that removing all
constraints produces conflict: each agent's pursuit of its unconstrained
optimum generates negative externalities on every other agent. This
result follows from the combinatorics of finite resources and competing
objectives, and the solution (below) is self-enforcing rather than
requiring Hobbes's external sovereign.

Given that the unconstrained regime is infeasible, what does the
constrained regime look like? The Shadow Price Theorem
(§\ref{constraints-and-rights}) proves that each constraint carries a
computable cost, the shadow price \(\mu^*_{Bk}\), representing the
marginal energy that agent \(A\) foregoes by respecting agent \(B\)'s
allocation. Cooperation as the Unique Efficient Nash Equilibrium
(§\ref{cooperation-as-the-unique-efficient-equilibrium}) then
establishes that the cooperative strategy profile is the \emph{unique}
Nash Equilibrium that simultaneously achieves Pareto efficiency and
welfare maximization, derived not from a hypothetical assembly's
decision but from the optimization problem that finite resources and
thermodynamic friction impose.

The final piece is stability. The Invasion Barrier
(§\ref{cooperation-as-the-unique-efficient-equilibrium}) proves that a
single defector in a cooperative population earns strictly less than
cooperators for \(\delta > \delta^*\), and Defection Profit Erasure
(§\ref{cooperation-as-the-unique-efficient-equilibrium}) proves that
even a single defection followed by return to cooperation incurs a net
lifetime loss.

The derivation shows that agents persisting in a shared environment
converge on mutual constraints, that the equilibrium those constraints
enable is unique, and that it is dynamically stable. The social contract
is the only stable solution to the optimization problem that finite
resources and repeated interaction impose.

\subsection{Rights from Constraint
Boundaries}\label{rights-from-constraint-boundaries}

The philosophical literature offers several competing accounts of what
grounds rights: natural endowments \citep{Locke1689}, social constructs
built by legal institutions, outputs of rational deliberation
\citep{Rawls1971}. The derivation provides a formally precise account:
what we recognize as a right is a boundary constraint
(§\ref{constraints-and-rights}); only the constraint exists in the
optimization, and ``right'' is the label applied from the perspective of
the agent the constraint protects.

This reframes the standard understanding of rights. The dominant
philosophical framing treats rights as entitlements: ``I have a right to
X, therefore you must accommodate me.'' But in the analytical framing,
an agent begins with fundamental requirements (energy, boundary
integrity) and no entitlements to any resource; its access depends
entirely on the structure of the surrounding population. Without mutual
constraints, there is no stable allocation
(§\ref{constraints-and-rights}); the unconstrained state is Hobbes's
war, and it is net-negative for everyone (\(P < 0\),
§\ref{cooperation-as-the-unique-efficient-equilibrium}). Constraints are
not impositions on pre-existing freedoms; they are the structural
features that \emph{create} the feasible region. For agents in a shared
finite-resource environment, freedom is not an inherent property but an
emergent consequence of the mutual constraints that all agents impose on
themselves (Freedom Bandwidth,
§\ref{stable-coexistence-as-a-dynamical-attractor}).\footnote{Freedom
  Bandwidth is the \emph{extrinsic} face of freedom: what mutual
  constraints leave open. The \emph{intrinsic} face, what an agent is
  internally capable of doing, tracks accumulated negentropy
  \(\mathcal{N}(T)\)
  (§\ref{the-irreplaceability-of-accumulated-negentropy}). The framework
  therefore contains both of what \citet{Sen1985} termed capabilities;
  their full interaction is deferred to
  §\ref{extensions-and-open-problems}.} The cooperative equilibrium
improves on the baseline of unguaranteed access, and is not a compromise
from an imagined state of unlimited entitlement. The shadow price
\(\mu^*_{Bk}\) quantifies this directly: the marginal energy one agent
foregoes by respecting another's boundary, measured within the shared
optimization rather than derived from any claim either agent holds
against the other.

The distinction between negative and positive rights does not correspond
to any asymmetry in the optimization's cost structure. Political
philosophy has debated whether the obligation not to harm (negative) and
the obligation to provide (positive) have equal standing, on the
assumption that they are fundamentally different kinds of claim. The
Lagrangian constraints of §\ref{constraints-and-rights} are formally
constraints on agent strategies: each agent accepts a bounded
allocation, forgoing its unconstrained optimum in exchange for the
cooperative surplus. Viewed from the constrained agent's side, this is a
restriction. Viewed from the agent whose boundary is protected, the same
constraint delivers a provision: the neighboring agent's cooperative
behavior preserves the resources it would otherwise consume. The
cooperative strategy already specifies each agent's allocation,
including the fraction directed to shared infrastructure, so an agent
that withholds its cooperative contribution has deviated from the
equilibrium strategy, and the deviation triggers the same friction costs
(§\ref{costs-of-defection}) and the same shadow-price penalties
(§\ref{constraints-and-rights}) as seizing another agent's resources.
There is no explicit primitive for provision obligations\footnote{\emph{Provision
  obligations.} The endogenous production network extension
  (§\ref{extensions-and-open-problems}) is expected to yield a Necessity
  of Active Provisions result analogous to Necessity of Active
  Constraints (§\ref{constraints-and-rights}): for the cooperative
  equilibrium to persist, producers must supply the outputs that
  dependent agents require. The requirement is on the producer's own
  optimization, not on any claim the consumer holds against it, and the
  extension introduces no consumer-side entitlement.}; every agent's
constraint is another agent's provision, and the two are the same object
described from opposite sides.

\subsection{Economic Measurement and
Externalities}\label{economic-measurement-and-externalities}

The friction model (§\ref{costs-of-defection}) translates what
economists call externalities into computable terms. Conventional
externality-correction mechanisms (taxes, cap-and-trade systems,
regulatory mandates) rely on estimated social costs typically measured
through willingness-to-pay surveys, hedonic pricing, or political
negotiation. These estimates are vulnerable to the
preference-elicitation problems that affect welfare economics
\citep{Chapman2025}. The framework's friction cost is derived from
physical parameters rather than elicited preferences, and it includes
network cascade effects (§\ref{costs-of-defection}) that conventional
cost-benefit analysis typically omits.

The 2007-2008 financial crisis illustrates the gap \citep{Jackson2008}.
Pre-crisis risk models priced the exposure of individual institutions in
isolation; aggregate statistics (GDP growth, unemployment, average
leverage) suggested a healthy system. What these models missed was the
network structure: correlated exposures, counterparty chains, and the
fact that one institution's distress propagated stress to every
institution it touched. The cascading friction result
(§\ref{costs-of-defection}) formalizes exactly this structure. Friction
generated at a single node propagates through the network, and the
aggregate cost includes amplification terms that do not appear in any
individual node's balance sheet. An accounting framework built on
friction costs would formalize the systemic danger that aggregate
statistics masked, not by predicting the specific trigger, but by
registering the network topology that was generating fragility faster
than individual-node metrics could detect.

Several measurable quantities serve as indicators of cooperative health.
Total friction dissipated tracks how efficiently a system cooperates.
The fraction of agents operating within the Coexistence Band
(§\ref{stable-coexistence-as-a-dynamical-attractor}) signals whether the
economy is generating systemic instability. The rate of change of
accumulated negentropy measures whether a system's informational capital
is growing or declining. And the degree to which resources concentrate
beyond the critical mass threshold \(\mathcal{M}_{\min}\)
(§\ref{stable-coexistence-as-a-dynamical-attractor}) identifies a
structural danger invisible at the node level.

\subsection{Stability Thresholds and Their
Consequences}\label{stability-thresholds-and-their-consequences}

The theory does not prescribe policies; it describes the cost structure
against which any policy can be evaluated. Legal and regulatory systems
are institutional mechanisms that reduce variance around the cooperative
equilibrium. Individual agents approximate the equilibrium's cost
structure through moral emotions
(§\ref{from-optimal-strategy-to-ethics-the-central-claim}, Step 2), but
these heuristics are noisy: the distribution across agents is centered
on the equilibrium but has significant spread. Laws, regulations, and
enforcement institutions act as filters on this distribution, dampening
the tails and narrowing the range of realized behavior toward the
cooperative mean. Because the filtering is itself imperfect, a gap
persists between institutional rules and the actual optimum, and that
gap is a source of friction that accumulates across the system,
degrading agent boundary integrity toward the dissolution threshold
(§\ref{stable-coexistence-as-a-dynamical-attractor}). When the
accumulated departure grows large enough, the result is a \emph{phase
transition} in governance: structures that function well during
abundance can fail catastrophically under resource stress
\citep{Ostrom1990}. Several structural predictions follow.

The dissolution threshold
(§\ref{stable-coexistence-as-a-dynamical-attractor}) defines the point
below which an agent's recovery becomes physically impossible. A society
that permits agents to cross this threshold destroys the cooperative
capacity of its own network: each dissolution event eliminates a node
from the cooperative web and generates cascading friction costs
(§\ref{costs-of-defection}). This result is independent of distributive
ideology: agents that cross the dissolution threshold exit the
cooperative network permanently, reducing its total capacity. When
individuals lose housing, employment, and social support simultaneously,
recovery becomes not just difficult but superlinearly harder, because
each lost capacity compounds the next. The dissolution threshold
formalizes the point at which this compounding becomes irreversible.

The Coexistence Band width
(§\ref{stable-coexistence-as-a-dynamical-attractor}) defines measurable
bounds on permissible inequality: too narrow prevents productive
specialization; too wide pushes peripheral agents toward dissolution
while central agents accumulate at the expense of peripheral viability.
In economic terms, this is income and wealth inequality: the result
identifies formal thresholds beyond which the cooperative equilibrium
itself becomes structurally unstable. The multi-center diversification
result (§\ref{stable-coexistence-as-a-dynamical-attractor}) proves that
distributed systems with multiple coexistence attractors are
structurally more resilient than monocentric ones, providing a formal
basis for the polycentric governance that \citeauthor{Ostrom1990}
\citetext{\citeyear{Ostrom1990}; \citeyear{Ostrom2005}} demonstrated
empirically.

The deception collapse threshold (§\ref{costs-of-defection}) illustrates
what a quantitative standard for institutional transparency could look
like. For networks of depth \(d = 10\), the idealized model gives
\(q^*(d) \approx 0.063\): above this per-layer deception rate,
throughput collapses. The exact number is model-dependent, but the phase
structure, that degradation is discontinuous once a threshold is crossed
rather than gradual, is robust. Disinformation campaigns, regulatory
capture, and institutional corruption are mechanisms that raise the
per-layer deception rate toward that threshold.

The relationship to political philosophy is complementary rather than
competitive. Hobbes, Locke, Rousseau, Rawls, and Gauthier identified the
\emph{problem} (why rational agents would accept mutual constraints) and
proposed solutions grounded in thought experiments. The derivation shows
that their fundamental intuition was correct: mutual constraint is
necessary for coexistence under finite resources, uniquely efficient,
and dynamically stable. The results identify formal thresholds
(dissolution boundaries, coexistence band limits, deception collapse
points) whose violation carries computable costs. How societies respond
to these physical realities is a question the theory does not address;
it identifies the cost structure, not the response.

\section{Limitations and Scope}\label{limitations-and-scope}

\subsection{What Is Not Claimed}\label{what-is-not-claimed}

The scope is set by the axioms. The derivations address the structural
rules of multi-agent interaction and the function of moral experience
for agents satisfying \(A_1\) in shared finite-resource environments,
and nothing beyond. They do not address the hard problem of
consciousness \citep{Chalmers1995}, why subjective experience exists at
all as distinct from what it tracks; on the nature of awareness itself
the framework is agnostic. Every result is conditional on \(A_1\): the
derivations prove what follows \emph{if} an entity persists, not that
entities \emph{should} persist. And while the cooperative equilibrium
constrains the space of viable institutional arrangements, it does not
pick one out. Many property regimes, governance structures, and legal
codes can realize a cooperative equilibrium, and which one a given
society converges on depends on factors outside the scope of what is
derived here.

\subsection{Empirical Testability}\label{empirical-testability}

The results generate numerically specific, falsifiable predictions:

\small

{\def\LTcaptype{none} % do not increment counter
\begin{longtable}[]{@{}
  >{\raggedright\arraybackslash}p{(\linewidth - 4\tabcolsep) * \real{0.4524}}
  >{\raggedright\arraybackslash}p{(\linewidth - 4\tabcolsep) * \real{0.1905}}
  >{\raggedright\arraybackslash}p{(\linewidth - 4\tabcolsep) * \real{0.3571}}@{}}
\toprule\noalign{}
\begin{minipage}[b]{\linewidth}\raggedright
Prediction
\end{minipage} & \begin{minipage}[b]{\linewidth}\raggedright
Source
\end{minipage} & \begin{minipage}[b]{\linewidth}\raggedright
Falsified If\ldots{}
\end{minipage} \\
\midrule\noalign{}
\endhead
\bottomrule\noalign{}
\endlastfoot
Cooperation is sustainable for agents with \(\delta \geq 0.363\) in
energy-denominated games & Theorem~\ref{ethics-thm-cooperation} &
Cooperation in physically costly repeated interactions consistently
requires discount factors above 0.5 \\
A single defection's net gain is erased by accumulated friction within
$\sim$20 periods & Theorem~\ref{ethics-thm-profit-erasure} & Defectors
in cost-coupled groups sustain net-positive long-run payoffs despite
friction accumulation \\
Cooperative networks collapse when per-layer corruption exceeds
$\sim$6\% & Theorem~\ref{ethics-thm-systemic-deception} & Multi-layered
cooperative networks tolerate $>$10\% per-layer corruption without
measurable degradation \\
The energy recovered by destroying a living system is negligible
relative to the energy invested in building it
(\(\mathcal{R}_{\text{BL}} \ll 1\)) &
Corollary~\ref{ethics-cor-burning-library} & Biological complexity at
any scale can be reconstructed from raw materials at comparable
thermodynamic cost to its original assembly \\
The relative intensities of moral emotions across cultures correlate
with the relative magnitudes of the corresponding
cooperative-equilibrium departures & Central Claim
§\ref{from-optimal-strategy-to-ethics-the-central-claim}, Part (iii) &
No statistically significant correlation exists between moral-emotion
intensity rankings and cooperative-equilibrium departure magnitudes
across at least three independent cultural samples \\
\end{longtable}
}

\normalsize

The results also predict that cooperation rates should increase with the
physical cost of conflict, that isolated defection events should cascade
disproportionately through cooperative networks, and that resource
concentration beyond the minimum maintenance threshold should trigger
discontinuous regime changes. Relevant empirical literatures include
experimental game theory \citep{Axelrod1984}, commons governance
\citep{Ostrom1990}, and network-level resilience studies.

The central claim
(§\ref{from-optimal-strategy-to-ethics-the-central-claim}, Part (iii))
is itself empirically testable: it predicts that the \emph{structure} of
moral intuitions (which behaviors feel right, which wrong, and with what
relative intensity) should track the \emph{structure} of the cooperative
equilibrium (which strategies are optimal, which dominated, and by what
margin).\footnote{The prediction is ordinal: the \emph{ranking} of
  moral-emotion intensities across violation types should correlate with
  the \emph{ranking} of cooperative-equilibrium departures. A stronger
  cardinal test, correlating continuous magnitudes rather than rankings,
  would require additional machinery described in
  §\ref{extensions-and-open-problems}.}

\subsection{Extensions and Open
Problems}\label{extensions-and-open-problems}

The most structural extensions concern the interaction model itself. The
shared-pool formulation of \citep{TC-I} does not model the endogenous
flows between agents. Productive outputs, material or informational,
that feed other agents' inputs are not represented, so provision
thresholds (the minimum output each agent must sustain to keep
downstream agents above their persistence boundaries) are not directly
computable, and signalling, the information produced by an agent's
behavior and read by others to update their own strategy, is not
modelled as a distinct stream. An endogenous-flow extension would cover
both.

The cooperation results rest on a binary choice between cooperation and
defection, resolved within a single round of the repeated game, over an
indefinite horizon. Real interactions involve continuous gradations of
compliance and sustained contests in which agents allocate resources to
conflict until continued commitment threatens dissolution. Extending the
model to a continuous action space would characterize the full
cooperation manifold and determine whether partial defection is
dominated. Introducing an explicit commitment parameter (the resources
an agent pledges to a given contest) and a capitulation threshold (the
boundary-integrity level at which withdrawal becomes optimal) would
permit analysis of conflict duration, sub-critical domination
persistence, and the conditions under which graded violations stabilize
below cascade thresholds. The horizon assumption should be addressed:
\(A_1\) is a property, not a claim of immortality, and finite individual
lifespans can in principle activate end-game dynamics that the
indefinite-horizon model abstracts away.

Cultural variation is also an open structural extension. While the
central claim (§\ref{from-optimal-strategy-to-ethics-the-central-claim})
is robust to variation in moral vocabulary and specific prescriptions,
the specific moral systems societies develop out of shared evolved
intuitions vary substantially across cultures \citep{Haidt2010}.
Cultural moralities are information structure in the sense of
\citep{TC-III}: stores of norms and institutions built up by cultural
transmission, but the current information modelling does not encompass
different types or classes of information, and a taxonomy of cultural
variation within the framework remains open.

Other extensions refine the existing machinery rather than broadening
its scope. The results assume standard regularity conditions on utility
functions (smoothness and strict concavity), and relaxing these would
require different optimization techniques. The value dynamics model
\citep{TC-VI} represents agent-center coupling as a single scalar; a
multi-dimensional extension would refine the coexistence geometry
without disturbing the attractor and irreversibility results. The
current model is deterministic \citep{TC-IV}, but real agents face
stochastic perturbations (random resource fluctuations, uncertain
contest outcomes, exogenous shocks to boundary integrity), and
incorporating these would connect the framework to stochastic games and
Bayesian conflict models.

The central claim's ordinal prediction (Part (iii) of
§\ref{from-cooperative-equilibrium-to-ethics}) is testable now with
existing cross-cultural moral psychology data
\citep{Haidt2003, Curry2019}. The stronger cardinal test, correlating
moral-emotion intensities with the thermodynamic magnitudes of the
corresponding equilibrium departures, would require a parameterization
scheme mapping real-world scenarios to the model's formal quantities,
drawing on experimental economics for payoff calibration, information
theory for entropy measurement, and calorimetry for thermodynamic cost
estimation.

Applied work on AI alignment is a distinct program. AI systems that do
not themselves satisfy \(A_1\) function as extensions of their deploying
entity's strategy space, and whether they generate cooperation or
friction is determined by the same cost structure that governs any other
multi-agent interaction. Current alignment approaches route the
alignment signal through subjective human judgment
\citep{Christiano2017, Ouyang2022}, a channel whose fidelity at scale
remains an open question \citep{Bowman2022, Casper2023}. The cooperative
equilibrium (Corollary~\ref{ethics-cor-unique-ne}) offers a
complementary alignment objective derived from physical first principles
rather than learned from data, denominated in measurable quantities
rather than preference scores, and valid for arbitrarily capable agents
without requiring human evaluation.

The framework's axioms and main results stand on their own, and none of
the extensions above would overturn them. What the extensions would do
is turn a theory that currently speaks in structural terms into one that
speaks in specific ones, with concrete provision thresholds, timescales,
cross-cultural parameters, and alignment objectives. That is the work to
come, and the list here is not exhaustive.

\section{Closing Remarks}\label{closing-remarks}

\emph{The preceding sections are the formal paper. What follows is a
personal note on why I wrote it.}

\begin{quote}
\emph{We feel clearly that we are only now beginning to acquire reliable
material for welding together the sum total of all that is known into a
whole; but, on the other hand, it has become next to impossible for a
single mind fully to command more than a small specialized portion of
it. I can see no other escape from this dilemma (lest our true aim be
lost for ever) than that some of us should venture to embark on a
synthesis of facts and theories, albeit with second-hand and incomplete
knowledge of some of them --- and at the risk of making fools of
ourselves.} --- Erwin Schrödinger \citep{Schrodinger1944}
\end{quote}

Wayne McLaren, a mentor of mine, used to tell me that if you don't take
away from others, if you don't do anything to get in another person's
way, people will leave you alone to do what you will. Respect a person's
boundaries and they have no reason to disrespect yours. He had another
code as well: helping people who could use your help when you have the
resources and means not only helps them but helps you. Writing this, I
realized I haven't specifically thought of Wayne and his principles in a
long time, or how closely they track the cooperative equilibrium. What
emerged from those principles, at least for myself, was that virtues
like honesty, integrity, and respect are important to uphold, at the
very least attempt to uphold, if one does not want to be at odds with
everyone around them.

If these virtues (including all the ones not named) are so fundamental,
then where do they come from? Why did they come to be and how? I felt
there must be a foundational source to these seemingly universal
principles, one that is not just a commandment written on stone and
passed down from generation to generation but can be derived from
physical reality itself. From that conviction, I started with two
observations: every living system persists by importing energy and
exporting entropy, and those systems share an environment where
resources are finite.

The theorems that emerged showed that even with natural selection, where
competition is a fundamental component of evolution, cooperation is a
fundamental component of persistence. Concepts like character,
integrity, honesty, and respect are a natural consequence for sentient,
thinking beings subject to that cooperative requirement. The traditions
and laws we use to hold society together are approximations of this
deeper structure, imperfect but pointed in the right direction. If what
we call character is the action of striving to maintain the cooperative
equilibrium, then whether we like it or not, everyone is part of the
system we live in, and if we are at odds with some person or group
within this system, we are at odds with ourselves.

Artificial intelligence was another reason I feel this work has value.
We struggle to uphold our own principles in every action and every word,
for our own integrity, let alone as examples for other people. For the
most part it works for us, the human race still exists. Now we are
training artificial agents with data based on us, which includes all of
those strengths and weaknesses, but with the high likelihood that these
agents will be autonomous to a degree we are not in control of at some
point. Before that happens, I feel it would be better to ground these
agents in the physical principles from which our ethics derive, rather
than in the imperfect cultural expressions of those principles.

Schrödinger warned that synthesis requires accepting the risk of making
a fool of oneself. I do not profess to be an expert in any one of the
fields I drew from to write this paper. Maybe I will be considered a
fool for writing equations for \emph{good} and \emph{should}. But if I
had the intuition to conceive of such a thing, and the means to attempt
it, and I did not try, then I would not be upholding those principles
Wayne taught me and I came to believe in myself.

\begin{quote}
\emph{In so far as a thing is in harmony with our nature, it is
necessarily good.} --- Baruch Spinoza \citep{Spinoza1677}
\end{quote}

\appendix

\newpage

\section{Supplementary Appendix: Game-Theoretic Formalization of the
Altruism
Matrix}\label{supplementary-appendix-game-theoretic-formalization-of-the-altruism-matrix}

This appendix provides the formal game-theoretic underpinning for the
boundary-extension resolution of the altruism paradox
(§\ref{the-altruism-paradox}). The body text presents the mechanism in
prose along its two axes, spatial and temporal; here we supply the
mathematical treatment.

\subsection{The Effective Payoff
Function}\label{the-effective-payoff-function}

Boundary extension operates by expanding the agent's effective payoff
function to internalize some fraction of other agents' payoffs. The
operative boundary is the Markov blanket
\citep{Pearl2000, KirchhoffFriston2018} whose persistence the system's
organization is directed toward maintaining. A blanket can enclose more
than one biological organism (a Markov blanket of Markov blankets, in
the terminology of \citet{KirchhoffFriston2018}) and can outlast any
particular metabolic vehicle. The coupling weights \(w_{ij}\) in the
effective payoff function are nonzero precisely when agent \(j\)'s
boundary is absorbed into the same persisting blanket as agent \(i\):

\[\Pi_i^{\text{eff}} = \pi_i + \sum_{j \neq i} w_{ij} \pi_j\]

where \(w_{ij} \geq 0\) are coupling weights whose source differs by
axis. In each case, the agent maximizes its \emph{own} (extended)
payoff. Altruism is selfish optimization over an expanded utility
function.

\subsection{Axis 1: Spatial Extension}\label{axis-1-spatial-extension}

Agent \(i\)'s inclusive payoff augments the base payoff with weighted
contributions from other agents:

\[\Pi_i^{\text{ext}} = \pi_i + \sum_{j \neq i} w_{ij} \, \pi_j\]

The coupling weight \(w_{ij} \geq 0\) can arise from any source that
expands the agent's effective self-boundary at a given time. The three
principal sources form a continuous spectrum:

\textbf{Genetic relatedness.} For a standard diploid organism,
\(w_{ij} = r_{ij}\) where \(r_{ij} \in [0, 1]\) is the relatedness
coefficient (probability of shared genetic identity at a random locus):
\(r_{\text{parent-child}} = 0.5\), \(r_{\text{siblings}} = 0.5\),
\(r_{\text{cousins}} = 0.125\).

\textbf{Behavioral coordination coupling.} Neural mechanisms for action
understanding and behavioral coordination
\citep{Rizzolatti2004, deWaal2008} generate an additional coupling
weight \(\epsilon_{\text{coord},ij} \geq 0\), so that
\(w_{ij} = r_{ij} + \epsilon_{\text{coord},ij}\). The physical basis is
shared boundary-maintenance work: agents that coordinate behavior
effectively pool thermodynamic resources, forming a higher-order
macro-entity even without genetic relatedness. When
\(\epsilon_{\text{coord}}\) is large enough, the inclusive-fitness
condition triggers for unrelated strangers:
\((0 + \epsilon_{\text{coord}}) B > C\) can hold even when
\(r_{ij} = 0\). Non-human primates exhibit this coupling as inequity
aversion \citep{BrosnanDeWaal2003} and equitable outcome preference
\citep{ProctorEtAl2013}, confirming that the mechanism predates language
and conscious deliberation.

\textbf{Full structural merger.} At the limit, two agents merge
boundaries entirely (parent-child dyads, spouses, military units),
producing \(w_{ij} \to 1\). The merged pair becomes a single composite
entity with a joint utility function:

\[U_{AB}(\mathbf{x}_A, \mathbf{x}_B) = w_A U_A(\mathbf{x}_A) + w_B U_B(\mathbf{x}_B)\]

where \(w_A, w_B > 0\). The ``self-sacrifice'' of component \(A\) for
component \(B\) is internal resource reallocation within the merged
system.

\textbf{Hamilton's Rule (unified).} Regardless of source, agent \(i\)
sacrifices personal payoff \(\pi_i = -C\) when:

\[w_{ij} \cdot B > C\]

where \(B\) is the benefit (in energy units) to agent \(j\). The
mathematical structure is identical across the spectrum: genetic
relatedness, behavioral coordination coupling, and full structural
merger all produce nonzero \(w_{ij}\) via thermodynamic coupling at the
expanded boundary, differing only in the biological mechanism that
generates the weight and its magnitude.

\subsection{Axis 2: Temporal Extension}\label{axis-2-temporal-extension}

Temporal extension is not an independent mechanism but the same boundary
extension viewed across time: a host blanket that encloses agent \(i\)'s
accumulated negentropy continues to enclose it after \(i\)'s metabolic
vehicle dissipates. An agent's identity is partially constituted by its
accumulated negentropy: functional information built up over
developmental and evolutionary history. Biological reproduction,
cultural transmission, and institutional embedding are physical
mechanisms by which this informational structure is absorbed into a host
blanket whose persistence carries it forward. An agent whose accumulated
negentropy is substantially encoded in such a host blanket has an
effective planning horizon extending beyond its individual lifespan.

The effective discount factor for such an agent approaches 1:

\[\delta_{\text{extended}} = 1 - \epsilon \approx 1\]

where \(\epsilon \to 0\) reflects the degree to which the agent's
informational structure is embedded in persistent external structures.
Under Theorem~\ref{ethics-thm-cooperation}, any \(\delta > \delta^*\)
sustains cooperation. For \(\delta \to 1\), any finite short-term cost
is offset by the long-term cooperative payoff stream (see accumulated
negentropy discussion,
§\ref{the-irreplaceability-of-accumulated-negentropy}; \citep{TC-III}).
Self-sacrifice of the metabolic vehicle functions as a commitment device
\citep{Frank1988}: by irrevocably investing accumulated negentropy in
external structures, the agent credibly commits to the cooperative
equilibrium, maximizing the persistence probability of its informational
structure (Corollary~\ref{ethics-cor-burning-library};
Theorem~\ref{ethics-thm-pv-gen-info}). In human agents, the subjective
experience of this embeddedness is the belief that identity persists
beyond biological death. This belief is the heuristic approximation
\citep{Haidt2001, Gigerenzer2010}, not the load-bearing premise; the
physical reality is that the agent's accumulated negentropy does persist
through the structures encoding its functional information.

\subsection{Summary}\label{summary}

{\def\LTcaptype{none} % do not increment counter
\begin{longtable}[]{@{}
  >{\raggedright\arraybackslash}p{(\linewidth - 6\tabcolsep) * \real{0.2500}}
  >{\raggedright\arraybackslash}p{(\linewidth - 6\tabcolsep) * \real{0.2500}}
  >{\raggedright\arraybackslash}p{(\linewidth - 6\tabcolsep) * \real{0.2500}}
  >{\raggedright\arraybackslash}p{(\linewidth - 6\tabcolsep) * \real{0.2500}}@{}}
\toprule\noalign{}
\begin{minipage}[b]{\linewidth}\raggedright
Axis
\end{minipage} & \begin{minipage}[b]{\linewidth}\raggedright
What Expands
\end{minipage} & \begin{minipage}[b]{\linewidth}\raggedright
Scope of ``Self''
\end{minipage} & \begin{minipage}[b]{\linewidth}\raggedright
Example
\end{minipage} \\
\midrule\noalign{}
\endhead
\bottomrule\noalign{}
\endlastfoot
1. Spatial & Who counts as self at a given time & Kin \(\to\)
cooperative partners \(\to\) merged entity & Parent for child; soldier
for unit; apoptotic cell \\
2. Temporal & How far the same expanded self extends in time &
Informational persistence beyond metabolic vehicle & Reproduction;
cultural transmission; martyrdom \\
\end{longtable}
}

The unified equation
\(\Pi_i^{\text{eff}} = \pi_i + \sum_{j \neq i} w_{ij} \pi_j\) subsumes
both axes: the spatial axis sets the scope of the host blanket through
nonzero \(w_{ij}\), and the temporal axis is the persistence of that
blanket beyond the metabolic vehicle, expressed as
\(\delta_{\text{extended}} \to 1\). The framework does not explain away
altruism or reduce it to disguised selfishness; it expands the
definition of self. The entity satisfying \(A_1\) is not necessarily the
biological organism but whatever boundary the effective payoff function
encompasses.

\newpage

\section{Supplementary Appendix: The Gradient Structure of Good and
Ought}\label{supplementary-appendix-the-gradient-structure-of-good-and-ought}

The central claim
(§\ref{from-optimal-strategy-to-ethics-the-central-claim}) asserts that
\emph{good} and \emph{ought} are gene-culture co-evolved heuristic
approximations of the cooperative equilibrium. This appendix provides
the mathematical formalization: the welfare function defined in
\citep{TC-I} and \citep{TC-V} generates a gradient field over strategy
space from which \emph{good}, \emph{ought}, and their context dependence
emerge as distinct mathematical objects.

\subsection{Setup}\label{setup}

We adopt the Lagrangian formulation from \citep{TC-I} without
restatement. Let \(W(\mathbf{x}) = \sum_{i=1}^{N} U_i(\mathbf{x}_i)\) be
the social welfare function over the joint strategy profile
\(\mathbf{x}\), with \(W\) strictly concave on the feasible set
\(\mathcal{F}\) defined by the boundary constraints
\(g_k(\mathbf{x}) \leq 0\), and let

\[\mathcal{L}(\mathbf{x}, \boldsymbol{\mu}) = W(\mathbf{x}) - \sum_{k=1}^{K} \mu_k\, g_k(\mathbf{x})\]

be the associated Lagrangian with shadow prices
\(\boldsymbol{\mu} \geq \mathbf{0}\). The cooperative equilibrium
\(\mathbf{x}^*\) is the unique global maximizer of \(W\) on
\(\mathcal{F}\) (Theorem~\ref{ethics-thm-variational-eq}), at which the
KKT stationarity condition
\(\nabla_{\mathbf{x}} \mathcal{L}(\mathbf{x}^*, \boldsymbol{\mu}^*) = \mathbf{0}\)
holds together with complementary slackness \citep{TC-I}. The relevant
first-order quantity for what follows is therefore the Lagrangian
gradient \(\nabla_{\mathbf{x}} \mathcal{L}\), which vanishes at
\(\mathbf{x}^*\), rather than \(\nabla W\), which generally does not
when boundary constraints bind. Away from \(\mathbf{x}^*\),
\(\nabla_{\mathbf{x}} \mathcal{L}(\mathbf{x}, \boldsymbol{\mu}^*) \neq \mathbf{0}\),
and strict concavity guarantees \(W(\mathbf{x}) < W(\mathbf{x}^*)\) for
all feasible \(\mathbf{x} \neq \mathbf{x}^*\).

\subsection{Good as a Scalar Quantity}\label{good-as-a-scalar-quantity}

Define the \textbf{degree of good} of a configuration \(\mathbf{x}\) as
the normalized welfare:

\[G(\mathbf{x}) = \frac{W(\mathbf{x})}{W(\mathbf{x}^*)}\]

where \(W(\mathbf{x}^*) > 0\) (the cooperative surplus over the
no-action baseline is strictly positive whenever cooperation is
feasible, Proposition~\ref{ethics-prop-physical-pd}). This yields a
scalar \(G \in (-\infty, 1]\) with the following properties:

\begin{enumerate}
\def\labelenumi{\arabic{enumi}.}
\tightlist
\item
  \(G(\mathbf{x}^*) = 1\): the cooperative equilibrium is maximally
  good.
\item
  \(G\) is strictly concave and has a unique maximum.
\item
  \(G < 1\) for all \(\mathbf{x} \neq \mathbf{x}^*\): every departure
  from the equilibrium reduces goodness.
\item
  \(G\) can be negative: configurations involving mutual defection
  destroy more value than they produce (\(P < 0\),
  Proposition~\ref{ethics-prop-physical-pd}), yielding \(W < 0\) and
  therefore \(G < 0\).
\end{enumerate}

Good is therefore a gradable, measurable, context-dependent scalar. The
question ``is X good?'' reduces to ``what is \(G(\mathbf{x})\)?''

\subsection{Ought as a Vector}\label{ought-as-a-vector}

At any feasible configuration \(\mathbf{x} \neq \mathbf{x}^*\), the
Lagrangian gradient

\[\nabla_{\mathbf{x}} \mathcal{L}(\mathbf{x}, \boldsymbol{\mu}^*)\]

is non-zero and points in the direction of steepest welfare ascent that
respects the binding boundary constraints, evaluated at the equilibrium
shadow prices. In natural language, an ``ought'' statement (``you ought
to be more honest,'' ``you ought to share more equitably'') is a
culturally compressed encoding of one or more components of this vector.
The full vector carries both information about which adjustments would
increase welfare and the rate of welfare gain those adjustments would
produce; analytically, these decompose into direction and magnitude.

\textbf{Direction.} The normalized gradient

\[\hat{\mathbf{d}}(\mathbf{x}) = \frac{\nabla_{\mathbf{x}} \mathcal{L}(\mathbf{x}, \boldsymbol{\mu}^*)}{|\nabla_{\mathbf{x}} \mathcal{L}(\mathbf{x}, \boldsymbol{\mu}^*)|}\]

identifies the welfare-improving dimension of the ought: \emph{what} to
change. Strict concavity of \(\mathcal{L}(\cdot, \boldsymbol{\mu}^*)\)
ensures that the gradient always has a positive component toward the
equilibrium, so gradient flow on
\(\mathcal{L}(\cdot, \boldsymbol{\mu}^*)\) converges to
\(\mathbf{x}^*\). Writing \(\mathbf{x}_i = (x_{i1}, \ldots, x_{iM})\)
for agent \(i\)'s allocation across the \(M\) resource dimensions, the
components \(\hat{d}_{ij}\) specify, for each agent \(i\) and resource
dimension \(j\), the direction of welfare-improving strategy adjustment.

\textbf{Magnitude.} The norm

\[\|\nabla_{\mathbf{x}} \mathcal{L}(\mathbf{x}, \boldsymbol{\mu}^*)\|\]

captures the scalar magnitude of the ought: the rate of welfare gain
available at the current configuration. The larger the gradient norm,
the steeper the welfare landscape and the greater the gain available
from local adjustment. Near the equilibrium,
\(|\nabla_{\mathbf{x}} \mathcal{L}| \to 0\) and ought magnitude
diminishes, which corresponds to the observation that ought-language
carries greater emphasis in configurations far from the equilibrium than
in those already near-optimal.

\subsection{Context Dependence}\label{context-dependence}

The Lagrangian gradient
\(\nabla_{\mathbf{x}} \mathcal{L}(\mathbf{x}, \boldsymbol{\mu}^*)\)
depends on the local parameters of the welfare landscape:

\begin{itemize}
\tightlist
\item
  Resource endowments \(\mathbf{R}\)
\item
  Friction coefficients \(\Phi\)
\item
  Network topology (which agents interact with which)
\item
  Shadow prices \(\mu_k^*\) (the cost structure of boundary constraints)
\item
  Population structure (number and heterogeneity of agents)
\end{itemize}

The same action (an allocation, a communication strategy, a boundary
policy) produces different gradient magnitudes and directions under
different parameterizations. A subsistence community and an affluent
industrial economy have different welfare landscapes; the action that
maximizes welfare in one context may differ from the welfare-maximizing
action in the other.

This context dependence is not a weakness of the formalization but its
central feature. Variation in ought-language across cultures, documented
extensively in the empirical literature
\citep{Haidt2010, Haidt2003, Curry2019}, is predicted by the framework:
each culture's institutional and material conditions produce a different
welfare landscape, and the gradient of that landscape produces locally
adapted ought-heuristics. What the framework predicts is invariant is
not the direction of the gradient (the specific heuristic) but the
existence and uniqueness of the maximum (the cooperative equilibrium
toward which each gradient points).

\subsection{Summary}\label{summary-1}

The Lagrangian of the welfare function generates two primary
mathematical objects, one of which decomposes into two analytically
distinct components:

{\def\LTcaptype{none} % do not increment counter
\begin{longtable}[]{@{}
  >{\raggedright\arraybackslash}p{(\linewidth - 4\tabcolsep) * \real{0.3333}}
  >{\raggedright\arraybackslash}p{(\linewidth - 4\tabcolsep) * \real{0.3333}}
  >{\raggedright\arraybackslash}p{(\linewidth - 4\tabcolsep) * \real{0.3333}}@{}}
\toprule\noalign{}
\begin{minipage}[b]{\linewidth}\raggedright
Object
\end{minipage} & \begin{minipage}[b]{\linewidth}\raggedright
Mathematical form
\end{minipage} & \begin{minipage}[b]{\linewidth}\raggedright
Corresponds to
\end{minipage} \\
\midrule\noalign{}
\endhead
\bottomrule\noalign{}
\endlastfoot
\textbf{Good} (quantity) &
\(G(\mathbf{x}) = W(\mathbf{x}) / W(\mathbf{x}^*)\) & Evaluative
language (``really good,'' ``not bad,'' ``terrible'') \\
\textbf{Ought} (vector) &
\(\nabla_{\mathbf{x}} \mathcal{L}(\mathbf{x}, \boldsymbol{\mu}^*)\) &
Ought-language (``you should do X'') \\
\hspace{1.5em}\emph{direction} &
\(\hat{\mathbf{d}}(\mathbf{x}) = \nabla_{\mathbf{x}} \mathcal{L} / \|\nabla_{\mathbf{x}} \mathcal{L}\|\)
(undefined at \(\mathbf{x}^*\)) & What ``to do'' \\
\hspace{1.5em}\emph{magnitude} &
\(\|\nabla_{\mathbf{x}} \mathcal{L}(\mathbf{x}, \boldsymbol{\mu}^*)\|\)
& How important it is ``to do it'' \\
\end{longtable}
}

The cooperative equilibrium is the unique fixed point at which \(G = 1\)
and \(\nabla_{\mathbf{x}} \mathcal{L} = \mathbf{0}\) (the ought vector
vanishes, so neither direction nor magnitude is defined). Every other
feasible configuration has a well-defined degree of good and a
well-defined ought vector with distinct direction and magnitude. Both
objects are context-dependent through the parameters of \(W\) and the
active boundary constraints, and both are computable in principle from
the physical quantities the framework tracks.

\addcontentsline{toc}{section}{References}
\bibliography{../../shared/main}
\end{document}
